\chapter{The divisor scheme and the representability of \(\mathrm{Pic}^0\)}
\label{chap:DivisorScheme}

Grothendieck's construction of the Picard scheme does not represent \(\mathrm{Pic}\)
directly. It represents first the functor of relative effective divisors, and then
presents \(\mathrm{Pic}\) as a quotient of that divisor scheme along the Abel map: the
Abel map is a surjection of sheaves whose fibres are linear systems, and a scheme
representing \(\mathrm{Pic}\) is obtained by descending along it
(\ref{def:localEquations_picClass}, and Kleiman's Theorem 4.8 and its proof).

This chapter builds the divisor half of that argument for the curve \(C/k\) of
Chapter~\ref{chap:Challenge}, in the form in which the representability seam
\ref{def:pic0RepresentableByOfCharts} can consume it. The three steps are:

\begin{enumerate}
  \item a functor \(\Div^{\,n}_{C/k}\) of degree-\(n\) relative divisors, defined on affine
    test algebras by \emph{local equations plus a finite chart adaptation carrying a
    colength certificate}, and its Zariski-localised form (Sections
    \ref{sec:divfam}--\ref{sec:divfamzar});
  \item an embedding \(\varepsilon\) of that functor into a product of two Grassmannians,
    sending a divisor \(D\) to the pair of window submodules
    \(H^0(\OO(aF - D)) \subseteq H^0(\OO(aF))\) at two pinned exponents \(a\), and a closed
    subscheme \(\mathbf{Div}\) of that Grassmannian pair carved out by the multiplication
    condition which the image satisfies (Section~\ref{sec:divscheme}); this is
    Grothendieck's embedding of a Quot-type functor into a Grassmannian by a twisted
    pushforward;
  \item the \emph{backward} half of representability --- every divisor class over an affine
    test induces a unique morphism to \(\mathbf{Div}\), and the class is recoverable from
    that morphism (Section~\ref{sec:divclassify}).
\end{enumerate}

The forward half --- a universal divisor family on \(\mathbf{Div}\) --- is not built here.
Section~\ref{sec:divrep} states a conditional package for it and isolates the unproved
statement (\ref{thm:divRepUniversalFamily}).  It does not by itself produce the
open-immersion certificates or the local-surjectivity instance consumed by
\ref{def:pic0RepresentableByOfCharts}; those are separate hypotheses at the representability
seam.

The finite two-chart presentation in this chapter is an auxiliary window engine.  The divisor
functor used in the representability argument instead permits certificates on arbitrary
affine-open covers of the relative curve.  Thus the fixed chart cover data below is not the
definition of that widened functor.

Throughout, \(k\) is a field, \(C\) is a smooth proper geometrically irreducible curve over
\(k\) of genus \(g\), \(\pi \colon C \to \PP^1_k\) is a finite dominant morphism, and
\(F\) is its fibre Weil divisor (\ref{def:fiberWeilDivisor}). For a \(k\)-algebra \(R\) we
write \(C_R\) for the relative curve (\ref{def:relCurve}) and \(V_0^R, V_1^R\) for the two
pinned charts of \(C_R\) obtained by pulling back the two standard charts of \(\PP^1\)
along \(\pi\) (\ref{def:relCover}, \ref{def:fiberTwoCover}). We abbreviate
\(H_a := H^0(C, \OO(aF))\) (\ref{def:divisorSections}).  These pinned charts occur only
in the auxiliary window calculations.  Kleiman's Main theorem, which is cited below, asks
that the morphism be projective (Zariski locally on the base), flat, and have integral
geometric fibres (\ref{thm:divRepUniversalFamily}); those hypotheses are not synonyms for
smooth, proper, and geometrically irreducible, and require an explicit bridge.

\section{Divisor equality of local-equation systems}
\label{sec:diveq}

A local-equation system (\ref{def:localEquations}) is a presentation of an effective
Cartier divisor, not the divisor itself: the same divisor is presented by many systems,
differing by refinement of the cover and by unit rescaling of the equations. The relation
below identifies exactly those, in one move.

\begin{lemma}[Restriction to a refinement]
  \label{lem:localEquations_restrict}
  \provenancetag{CUSTOM}
  \lean{AlgebraicGeometry.Scheme.LocalEquations.restrict,
    AlgebraicGeometry.Scheme.LocalEquations.picClass_restrict}
  \uses{def:localEquations, def:localEquations_picClass, lem:localEquations_ratioUnit}
  \leanok
  Let \(d\) be a local-equation system on \(X\) adapted to a pointed cover \(\mathcal U\),
  and let \(\mathcal V \le \mathcal U\) be a refinement (\(\mathcal V(x) \subseteq
  \mathcal U(x)\) for every \(x\)). Restricting each equation, \(f'_x := f_x|_{\mathcal
  V(x)}\), is again a local-equation system, and it has the same Picard class:
  \(\mathrm{picClass}(d|_{\mathcal V}) = \mathrm{picClass}(d)\).
\end{lemma}
\begin{proof}
  \leanok
  Restriction of sections is a ring homomorphism and germs are unchanged along it, so each
  \(f'_x\) is again regular at every point of \(\mathcal V(x)\); restricting the overlap
  equation \(f_x|_\cap = g_{xy} f_y|_\cap\) of \(d\) (\ref{lem:localEquations_ratioUnit})
  to \(\mathcal V(x) \cap \mathcal V(y)\) shows the overlap condition, with transition unit
  the restriction \(g_{xy}|_\cap\). Thus the unit cocycle of \(d|_{\mathcal V}\) is the
  restriction of the cocycle of \(d\) along the refinement \(\mathcal V \le \mathcal U\).
  In the colimit defining \(\Pic(X)\) (\ref{def:localEquations_picClass}) a cocycle and its
  restriction along a refinement have the same class, by the definition of the colimit.
\end{proof}

\begin{lemma}[Divisor equality]
  \label{lem:localEquations_divEq}
  \provenancetag{CUSTOM}
  \lean{AlgebraicGeometry.Scheme.LocalEquations.DivEq,
    AlgebraicGeometry.Scheme.LocalEquations.divEq_refl,
    AlgebraicGeometry.Scheme.LocalEquations.DivEq.symm,
    AlgebraicGeometry.Scheme.LocalEquations.DivEq.trans,
    AlgebraicGeometry.Scheme.LocalEquations.picClass_eq_of_divEq}
  \uses{def:localEquations, lem:localEquations_restrict, lem:localEquations_rescale,
    def:localEquations_picClass}
  \leanok
  \source{kleiman-picard}
  Call two local-equation systems \(d_1, d_2\) on \(X\) \emph{divisor-equal}, written
  \(d_1 \sim d_2\), if there is a pointed cover \(\mathcal W\) refining both covers such
  that for every \(x \in X\) the two restricted equations differ by a unit on
  \(\mathcal W(x)\):
  \[
    f^{(1)}_x\big|_{\mathcal W(x)} \;=\; u_x \cdot f^{(2)}_x\big|_{\mathcal W(x)},
    \qquad u_x \in \struct X(\mathcal W(x))^{\times}.
  \]
  This is an equivalence relation, and \(\mathrm{picClass}\) is constant on its classes:
  \(d_1 \sim d_2\) implies \(\struct X(D_1) = \struct X(D_2)\) in \(\Pic(X)\).
\end{lemma}
\begin{proof}
  \leanok
  \emph{Reflexivity}: take \(\mathcal W = \mathcal U\) and \(u_x = 1\).
  \emph{Symmetry}: the same \(\mathcal W\) with \(u_x^{-1}\).
  \emph{Transitivity}: given \(d_1 \sim d_2\) on \(\mathcal W\) with units \(u\), and
  \(d_2 \sim d_3\) on \(\mathcal V\) with units \(v\), take the pointwise intersection
  \(\mathcal W \wedge \mathcal V\), which refines all three covers. Restricting the two
  equations to it and multiplying,
  \[
    f^{(1)}_x = (u_x|)\,f^{(2)}_x = (u_x|)(v_x|)\, f^{(3)}_x ,
  \]
  and \((u_x|)(v_x|)\) is a unit, being a product of restrictions of units.

  \emph{Invariance of the class}: restrict both systems to the witnessing cover
  \(\mathcal W\); by Lemma~\ref{lem:localEquations_restrict} neither class moves. On
  \(\mathcal W\) the first restricted system is the second one rescaled by the unit family
  \((u_x)_x\), so the two classes agree by Lemma~\ref{lem:localEquations_rescale}.
\end{proof}

\section{The relative divisor functor of the curve}
\label{sec:divfam}

Fix a \(k\)-algebra \(R\). A relative effective divisor on \(C_R/R\) is, locally on
\(C_R\), cut out by one equation which is regular on the fibres; a local-equation system on
\(C_R\) records exactly that, but on an \emph{arbitrary} pointed cover, and with regularity
asked in the stalks of \(C_R\) rather than on the fibres.  For the auxiliary pinned-window
presentation, a refinement by finitely many basic opens of the two pinned charts is paired
with a certificate asserting that the resulting colength modules are finite projective over
\(R\) of constant rank \(n\).  The widened divisor carrier used by the representability route
is a separate quotient in which each local certificate is carried by an arbitrary affine-open
cover of the relative curve.  The \(\Div^{\,n}_{\mathrm{Zar}}\) carrier defined below retains
the pinned-window certificate and is used only for the auxiliary comparison.
In either presentation the certificate is the proposed formal
analogue of \(R\)-flatness with fibres of degree \(n\) in Kleiman's Definition of
\(\Div^{\,n}_{X/S}\).

\begin{definition}[Auxiliary finite chart cover data]
  \label{def:finCoverData}
  \provenancetag{CUSTOM}
  \lean{AlgebraicGeometry.FinCoverData, AlgebraicGeometry.FinCoverData.pieces,
    AlgebraicGeometry.FinCoverData.cover₀, AlgebraicGeometry.FinCoverData.cover₁,
    AlgebraicGeometry.FinCoverData.isAffineOpen_pieces,
    AlgebraicGeometry.le_iSup_basicOpen_of_sum_eq_one}
  \uses{def:relCurve, def:relCover, def:fiberTwoCover}
  \leanok
  This auxiliary pinned-window datum on \(C_R\) consists of finitely many sections
  \(h^0_1,\dots,h^0_{m_0} \in \struct{}(V_0^R)\) and \(h^1_1,\dots,h^1_{m_1} \in
  \struct{}(V_1^R)\), together with coefficients \(a^\iota_j\) in the same rings satisfying
  the two partition identities
  \[
    \sum_{j} a^0_j h^0_j = 1 \ \text{ in } \struct{}(V_0^R), \qquad
    \sum_{j} a^1_j h^1_j = 1 \ \text{ in } \struct{}(V_1^R).
  \]
  Its \emph{pieces} are the basic opens \(P_j := D(h^\iota_j) \subseteq V_\iota^R\), indexed
  by the disjoint union of the two index sets. Each piece is affine, and the pieces of each
  chart cover that chart:
  \(V_0^R = \bigcup_j D(h^0_j)\) and \(V_1^R = \bigcup_j D(h^1_j)\).
\end{definition}
\begin{proof}
  \leanok
  A basic open of an affine open is affine, and the pinned charts are affine over the
  base, so each \(P_j\) is affine.

  For the covering claim on \(V_0^R\): let \(z \in V_0^R\) and suppose \(z \notin D(h^0_j)\)
  for every \(j\), that is, the germ of every \(h^0_j\) at \(z\) lies in the maximal ideal
  of \(\struct{C_R, z}\). Taking germs at \(z\) in the partition identity gives
  \(1 = \sum_j (a^0_j)_z (h^0_j)_z\), which exhibits \(1\) as an element of the maximal
  ideal --- impossible in a local ring. Hence some \(h^0_j\) is a unit at \(z\), i.e.\
  \(z \in D(h^0_j)\). The same argument applies on \(V_1^R\).
\end{proof}

\begin{definition}[Finite chart adaptation of a local-equation system]
  \label{def:divisorAdaptation}
  \provenancetag{CUSTOM}
  \lean{AlgebraicGeometry.DivisorAdaptation, AlgebraicGeometry.DivisorAdaptation.ofAnchors}
  \uses{def:finCoverData, def:localEquations, lem:localEquations_ratioUnit}
  \leanok
  \source{kleiman-picard}
  Let \(d\) be a local-equation system on \(C_R\), adapted to a pointed cover
  \(\mathcal U\). A \emph{finite chart adaptation} \(A\) of \(d\) consists of the auxiliary
  finite chart cover data (\ref{def:finCoverData}) with pieces \(P_j\), together with sections
  \(f_j \in \struct{}(P_j)\) subject to the \emph{refinement clause}: for every piece index
  \(j\) and every point \(y \in C_R\) there is a unit
  \(u \in \struct{}(P_j \cap \mathcal U(y))^{\times}\) with
  \[
    f_j\big|_{P_j \cap \mathcal U(y)} \;=\; u \cdot d_y\big|_{P_j \cap \mathcal U(y)} .
  \]
  (The quantifier ranges over \emph{all} \(y\); over an empty overlap the section ring is
  the zero ring and the clause is vacuous.)

  In particular, anchored data suffice to build an adaptation: if for each \(j\) one is
  given a point \(p_j\) with \(P_j \subseteq \mathcal U(p_j)\) and a unit \(w_j\) on
  \(P_j\) with \(f_j = w_j \cdot d_{p_j}|_{P_j}\), then the refinement clause holds for
  every \(y\), with unit the product of \(w_j|\) and the restriction of the transition unit
  \(g_{p_j y}\) of \(d\) (\ref{lem:localEquations_ratioUnit}).
\end{definition}
\begin{proof}
  \leanok
  Only the last assertion needs an argument. Fix \(j\) and \(y\), and restrict everything
  to \(W := P_j \cap \mathcal U(y)\). Restricting \(f_j = w_j\, d_{p_j}|_{P_j}\) to \(W\)
  gives \(f_j|_W = (w_j|_W)\, d_{p_j}|_W\). Restricting the defining equation of the
  transition unit, \(d_{p_j}|_{\mathcal U(p_j) \cap \mathcal U(y)} = g_{p_j y}\,
  d_y|_{\mathcal U(p_j)\cap\mathcal U(y)}\), along \(W \subseteq \mathcal U(p_j) \cap
  \mathcal U(y)\) gives \(d_{p_j}|_W = (g_{p_j y}|_W)\, d_y|_W\). Substituting,
  \(f_j|_W = (w_j|_W)(g_{p_j y}|_W)\, d_y|_W\), and the left factor is a unit on \(W\), a
  product of restrictions of units.
\end{proof}

\begin{lemma}[The adapted equations are regular]
  \label{lem:divisorAdaptation_eqn_regular}
  \provenancetag{CUSTOM}
  \lean{AlgebraicGeometry.DivisorAdaptation.eqn_regular}
  \uses{def:divisorAdaptation, def:localEquations}
  \leanok
  Let \(A\) be a finite chart adaptation of \(d\). For every piece index \(j\) and every
  point \(z \in P_j\), the germ of \(f_j\) at \(z\) is a nonzerodivisor in
  \(\struct{C_R,z}\).
\end{lemma}
\begin{proof}
  \leanok
  Apply the refinement clause of \ref{def:divisorAdaptation} at the piece \(j\) and at the
  point \(y := z\) itself. The point \(z\) lies in \(P_j\) and in \(\mathcal U(z)\), hence
  in the overlap \(W := P_j \cap \mathcal U(z)\), so germs at \(z\) may be taken in the
  identity \(f_j|_W = u\, d_z|_W\); germs commute with restriction, so
  \[
    (f_j)_z \;=\; u_z \cdot (d_z)_z \quad\text{in } \struct{C_R,z}.
  \]
  The germ \(u_z\) of a unit section is a unit, hence a nonzerodivisor; and \((d_z)_z\) is
  a nonzerodivisor by the regularity clause of \(d\) (\ref{def:localEquations}). A product
  of nonzerodivisors is a nonzerodivisor.
\end{proof}

\begin{definition}[The chart colength modules and their glued equalizer]
  \label{def:gluedColength}
  \provenancetag{CUSTOM}
  \lean{AlgebraicGeometry.DivisorAdaptation.colength,
    AlgebraicGeometry.DivisorAdaptation.ovlColength,
    AlgebraicGeometry.DivisorAdaptation.gluedSubmodule,
    AlgebraicGeometry.DivisorAdaptation.mem_gluedSubmodule_iff}
  \uses{def:divisorAdaptation, lem:divisorAdaptation_eqn_regular}
  \leanok
  Let \(A\) be a finite chart adaptation of \(d\), with pieces \(P_j\) and equations
  \(f_j\). Its \emph{chart colength modules} are the \(R\)-algebras
  \[
    Q_j \;:=\; \struct{}(P_j) \big/ (f_j),
  \]
  the structure sheaves of the pieces of the divisor cut out by \(d\). On an overlap put
  \(Q_{ij} := \struct{}(P_i \cap P_j)/(f_i|, f_j|)\), the quotient by the ideal generated
  by \emph{both} restricted equations, and let \(\delta^-, \delta^+ \colon \prod_j Q_j \to
  \prod_{(i,j)} Q_{ij}\) be the two restriction maps, \(\delta^-\) restricting the
  \(i\)-component and \(\delta^+\) the \(j\)-component. The \emph{glued colength module} is
  their equalizer
  \[
    W(A) \;:=\; \ker\bigl(\delta^- - \delta^+\bigr) \;\subseteq\; \prod_j Q_j ,
  \]
  an \(R\)-submodule; concretely \(s \in W(A)\) iff \(s_i\) and \(s_j\) have the same image
  in \(Q_{ij}\) for every pair \((i,j)\). It is the module of global sections of the
  structure sheaf of the divisor over the union of the pieces.
\end{definition}

\begin{definition}[The colength certificate of degree \(n\)]
  \label{def:isCertified}
  \provenancetag{CUSTOM}
  \lean{AlgebraicGeometry.DivisorAdaptation.IsCertified}
  \uses{def:gluedColength}
  \leanok
  \source{kleiman-picard}
  A finite chart adaptation \(A\) is \emph{certified of degree \(n\)} if
  \begin{enumerate}
    \item[(c1)] every chart colength module \(Q_j\) is finite and projective over \(R\);
    \item[(c2)] the glued colength module \(W(A)\) is finite and projective over \(R\), and
      has constant rank \(n\) at every prime: \(\operatorname{rank}_{\mathfrak p} W(A) = n\)
      for all \(\mathfrak p \in \Spec R\);
    \item[(c3)] the cokernel of the inclusion \(W(A) \hookrightarrow \prod_j Q_j\) is flat
      over \(R\);
    \item[(c4)] the cokernel of \(\delta^- - \delta^+\) is flat over \(R\).
  \end{enumerate}
  Clauses (c1) and (c2) say that the divisor is finite flat of degree \(n\) over the base
  --- Kleiman's condition that a relative effective divisor of degree \(n\) have
  \(S\)-flat structure sheaf with fibres of length \(n\). Clauses (c3) and (c4) are what
  make the certificate stable under an \emph{arbitrary} base change \(R \to R'\): they are
  the flatness hypotheses under which tensoring the defining left-exact sequence
  \(0 \to W(A) \to \prod_j Q_j \to \prod_{ij} Q_{ij}\) with \(R'\) stays exact, so that the
  base-changed glued module is again the equalizer of the base-changed arrows. Over a field
  they are automatic.
\end{definition}

\begin{definition}[Certified divisor families and the divisor functor value]
  \label{def:certifiedDivisorFamily}
  \provenancetag{CUSTOM}
  \lean{AlgebraicGeometry.CertifiedDivisorFamily, AlgebraicGeometry.divFamSetoid,
    AlgebraicGeometry.DivFam, AlgebraicGeometry.DivFam.mk,
    AlgebraicGeometry.DivFam.mk_eq_mk_iff, AlgebraicGeometry.DivFam.picClass}
  \uses{def:isCertified, lem:localEquations_divEq, def:localEquations_picClass}
  \leanok
  \source{kleiman-picard}
  A \emph{certified divisor family of degree \(n\)} over \(R\) is a triple \((d, A, c)\):
  a local-equation system \(d\) on \(C_R\), a finite chart adaptation \(A\) of \(d\), and a
  certificate \(c\) of degree \(n\) for \(A\). Divisor equality of the underlying systems
  (\ref{lem:localEquations_divEq}) is an equivalence relation on such triples --- the
  adaptation and the certificate are presentation data and do not enter the relation ---
  and
  \[
    \Div^{\,n}(R) \;:=\; \bigl\{\,\text{certified divisor families of degree }n\text{ over }
      R\,\bigr\} \big/ \sim
  \]
  is the value of this certified subfunctor at \(R\). The Picard class descends to it:
  \(\mathcal{O}(D) \in \Pic(C_R)\) is well defined on \(\Div^{\,n}(R)\).  On a certified
  triple this is the Abel-class map; it agrees with Kleiman's Abel construction only after
  the triple is shown to represent a flat effective divisor of degree \(n\).  No theorem here
  says that every effective flat Cartier divisor admits this pinned certificate.
\end{definition}
\begin{proof}
  \leanok
  The relation is an equivalence relation because divisor equality is
  (\ref{lem:localEquations_divEq}), and it ignores the last two components by definition.
  The Picard class of a certified family depends only on \(d\), and is constant on
  divisor-equality classes by the invariance half of \ref{lem:localEquations_divEq}; hence
  it factors through the quotient.
\end{proof}

\section{The Zariski-local certificate}
\label{sec:divfamzar}

The certificate of \ref{def:isCertified} is a condition on \emph{global} modules over
\(R\), and it is genuinely restrictive: it forces the divisor to sit, chart-trace by
chart-trace, in a closed way inside the two pinned charts.  The following is the obstruction
for this auxiliary pinned certificate.  It is not a no-go theorem for the widened functor,
whose pieces may be arbitrary affine opens, and it is the reason that carrier is kept
separate from this window engine.

\begin{theorem}[The chart obstruction to a global certificate]
  \label{thm:isCertified_chart_obstruction}
  \provenancetag{CUSTOM}
  \lean{AlgebraicGeometry.DivisorAdaptation.supportLeak_eq_empty_iff_finite_colength,
    AlgebraicGeometry.DivisorAdaptation.isClosed_supportLocus_inter_chart_of_isCertified,
    AlgebraicGeometry.DivisorAdaptation.not_isCertified_of_not_isClosed_inter_chart₀,
    AlgebraicGeometry.DivisorAdaptation.not_isCertified_of_divEq_of_isPreconnected_of_witnesses}
  \uses{def:isCertified, lem:localEquations_divEq, def:relCover}
  \leanok
  Write \(\operatorname{Supp}(d) \subseteq C_R\) for the support of \(d\), the closed set
  of points at which the equation of \(d\) is not a unit.
  \begin{enumerate}
    \item At a single piece \(P_j\), the chart colength module \(Q_j\) is a finite
      \(R\)-module if and only if \(\operatorname{Supp}(d) \cap P_j\) is closed in \(C_R\)
      --- that is, iff the support does not \emph{leak} out of the piece.
    \item Consequently, if \(A\) is a certified adaptation of \(d\) in any degree, then both
      pinned-chart traces \(\operatorname{Supp}(d) \cap V_0^R\) and
      \(\operatorname{Supp}(d) \cap V_1^R\) are closed in \(C_R\).
    \item In particular a divisor one of whose chart traces is not closed admits no
      certified adaptation in any degree; and if some representative of a divisor class has
      preconnected support containing a point outside \(V_0^R\) and a point outside
      \(V_1^R\), then \emph{no} representative of that class admits one.
  \end{enumerate}
\end{theorem}
\begin{proof}
  \leanok
  Say the \emph{leak} of the piece \(P_j\) is
  \(\overline{\operatorname{Supp}(d) \cap P_j} \setminus P_j\); it is empty exactly when
  \(\operatorname{Supp}(d) \cap P_j\) is closed in \(C_R\) (a set is closed iff it contains
  its closure, and the trace is closed in the open \(P_j\) always).

  (i) The piece \(P_j\) is affine (\ref{def:finCoverData}), and
  \(\operatorname{Supp}(d) \cap P_j\) is the closed subscheme \(Z_j := V(f_j) \subseteq
  P_j\) with coordinate ring \(Q_j\) --- a point of \(P_j\) is a support point exactly when
  \(f_j\) is not a unit there, i.e.\ when \(f_j\) lies in the corresponding prime.

  If \(Q_j\) is a finite \(R\)-module, then \(Z_j \to \Spec R\) is a finite, hence proper,
  morphism. The structure morphism \(C_R \to \Spec R\) is separated (it is a base change of
  the proper \(C \to \Spec k\)), so the immersion \(Z_j \to C_R\) over \(\Spec R\) is
  proper, and the image of a proper morphism is closed. Hence
  \(\operatorname{Supp}(d) \cap P_j\) is closed in \(C_R\), i.e.\ the leak is empty.

  Conversely, if \(\operatorname{Supp}(d) \cap P_j\) is closed in \(C_R\), then \(Z_j\) is a
  closed subscheme of \(C_R\) contained in the affine open \(P_j\); \(C_R \to \Spec R\) is
  proper, so \(Z_j \to \Spec R\) is proper, and a proper affine morphism is finite. Hence
  \(Q_j\) is a finite \(R\)-module.

  (ii) Clause (c1) of the certificate supplies finiteness of \(Q_j\) at every piece, hence
  by (i) an empty leak at every piece. The pieces indexed by the first chart cover
  \(V_0^R\) (\ref{def:finCoverData}), so \(\operatorname{Supp}(d) \cap V_0^R\) is the finite
  union \(\bigcup_j (\operatorname{Supp}(d) \cap P_j)\) of sets closed in \(C_R\), hence
  closed in \(C_R\). Likewise for \(V_1^R\).

  (iii) The first assertion is the contrapositive of (ii). For the second, first note that
  divisor equality does not move the support: on a common refinement the two equations
  differ by a unit at every point, and being a unit is a germ-local condition, so the two
  systems have the same non-unit locus. Hence the connectivity hypothesis and the two
  witnesses transport from the given representative to any other, and it suffices to rule
  out a certified adaptation of a system \(d\) whose own support is preconnected and meets
  the complements of both charts.

  Suppose it had one. By (ii) both traces \(\operatorname{Supp}(d) \cap V_\iota^R\) are
  closed in \(C_R\), hence closed in \(\operatorname{Supp}(d)\); each is also open in
  \(\operatorname{Supp}(d)\), being the intersection with an open. So each trace is clopen
  in the preconnected space \(\operatorname{Supp}(d)\), hence is empty or everything. If
  \(\operatorname{Supp}(d) \cap V_0^R = \operatorname{Supp}(d)\) then
  \(\operatorname{Supp}(d) \subseteq V_0^R\), contradicting the witness outside \(V_0^R\);
  otherwise \(\operatorname{Supp}(d) \cap V_0^R = \emptyset\), and since
  \(V_0^R \cup V_1^R = C_R\) this forces \(\operatorname{Supp}(d) \subseteq V_1^R\),
  contradicting the witness outside \(V_1^R\).
\end{proof}

\begin{definition}[The Zariski-local certificate]
  \label{def:isLocallyCertified}
  \provenancetag{CUSTOM}
  \lean{AlgebraicGeometry.IsLocallyCertified, AlgebraicGeometry.IsLocallyCertified.of_divEq,
    AlgebraicGeometry.CertifiedDivisorFamily.isLocallyCertified}
  \uses{def:certifiedDivisorFamily, def:relCurveMap, def:localEquations_pullback,
    lem:localEquations_divEq}
  \leanok
  A local-equation system \(d\) on \(C_R\) is \emph{locally certified of degree \(n\)} if
  there are finitely many \(g_1, \dots, g_m \in R\) generating the unit ideal such that for
  each \(i\) there is a certified divisor family of degree \(n\) over the localisation
  \(R_{g_i}\) whose local-equation system is divisor-equal to the pullback of \(d\) along
  \(C_{R_{g_i}} \to C_R\).

  Being locally certified is invariant under divisor equality, and every certified family
  is locally certified.
\end{definition}
\begin{proof}
  \leanok
  The comparison \(C_{R_{g}} \to C_R\) of an away localisation is an open immersion, so
  pullback of a local-equation system along it is defined without further hypothesis: the
  regularity required by \ref{def:localEquations_pullback} holds because an open immersion
  induces isomorphisms on stalks, and an isomorphism carries nonzerodivisors to
  nonzerodivisors.

  \emph{Invariance}: if \(d \sim d'\) then their pullbacks along each \(C_{R_{g_i}} \to
  C_R\) are divisor-equal (restrict the witnessing cover and pull the comparison units
  back; pullback of a unit is a unit), so the recorded divisor equalities for \(d\) compose
  with these to give the ones needed for \(d'\) (\ref{lem:localEquations_divEq},
  transitivity).

  \emph{Global implies local}: take \(m = 1\) and \(g_1 = 1\), whose span is the unit
  ideal. Over \(R_1\) take the base change of the given certified family; the certificate
  base-changes by clauses (c3)/(c4) of \ref{def:isCertified}, and its local-equation system
  is the pullback of \(d\) on the nose, so the required divisor equality is reflexivity.
\end{proof}

\begin{theorem}[Regularity of the pulled equations of a locally certified system]
  \label{thm:isLocallyCertified_pullback_regular}
  \provenancetag{CUSTOM}
  \lean{AlgebraicGeometry.IsLocallyCertified.germ_pullbackEqn_mem_nonZeroDivisors}
  \uses{def:isLocallyCertified, def:localEquations_pullback, def:relCurveMap}
  \leanok
  Let \(d\) be locally certified of degree \(n\) over \(R\), and let \(R \to R'\) be an
  arbitrary morphism of \(k\)-algebras. Then the equations of \(d\) pulled back along
  \(C_{R'} \to C_R\) are regular; hence the pullback \(d|_{R'}\) is again a local-equation
  system (\ref{def:localEquations_pullback}), with no flatness hypothesis on \(R \to R'\).
\end{theorem}
\begin{proof}
  \leanok
  Regularity of a section is a condition on its germ at each point, hence may be checked on
  any open cover of \(C_{R'}\). Let \(g_1,\dots,g_m\) be the unit-ideal family certifying
  \(d\). Their images \(\varphi(g_i) \in R'\) again generate the unit ideal, so the open
  immersions \(C_{R'_{\varphi(g_i)}} \to C_{R'}\) cover \(C_{R'}\); it suffices to prove
  regularity of the pulled germs at points of each such chart.

  Fix \(i\). The localisation map factors: \(R \to R_{g_i} \to R'_{\varphi(g_i)}\), the
  second arrow being the induced map of away localisations. Correspondingly
  \[
    C_{R'_{\varphi(g_i)}} \longrightarrow C_{R_{g_i}} \longrightarrow C_R
  \]
  composes to the restriction of \(C_{R'} \to C_R\), so the pulled equations over
  \(R'_{\varphi(g_i)}\) are the equations of \(d\) pulled first to \(C_{R_{g_i}}\) and then
  along the second arrow. On \(C_{R_{g_i}}\) the pulled system is divisor-equal to the
  local-equation system of a \emph{certified} family \(G_i\) by hypothesis. For a certified
  family the pulled equations along an arbitrary base change are regular: clause (c1) of
  the certificate makes each chart colength module projective, so the multiplication-by-
  \(f_j\) sequence \(0 \to \struct{}(P_j) \to \struct{}(P_j) \to Q_j \to 0\) is split as
  \(R\)-modules and stays exact after \(- \otimes_R R'\), which is precisely the statement
  that the pulled equation is again a nonzerodivisor. Finally, divisor equality is
  preserved by pullback and units pull back to units, so regularity transports from the
  pulled equations of \(G_i\) to those of \(d\).
\end{proof}

\begin{definition}[The auxiliary locally certified divisor functor \(\Div^{\,n}_{\mathrm{Zar}}\)]
  \label{def:divFamZar}
  \provenancetag{CUSTOM}
  \lean{AlgebraicGeometry.DivFamZar, AlgebraicGeometry.DivFamZar.mk,
    AlgebraicGeometry.DivFamZar.mk_eq_mk_iff, AlgebraicGeometry.DivFamZar.picClass,
    AlgebraicGeometry.DivFam.toZar}
  \uses{def:isLocallyCertified, lem:localEquations_divEq, def:certifiedDivisorFamily}
  \leanok
  The value at \(R\) of the \emph{locally certified} divisor functor is
  \[
    \Div^{\,n}_{\mathrm{Zar}}(R) \;:=\;
      \bigl\{\, d \text{ locally certified of degree } n \text{ on } C_R \,\bigr\} \big/ \sim ,
  \]
  the quotient by divisor equality (well posed by \ref{def:isLocallyCertified}). The Picard
  class descends to it, and passing from a global to a local certificate gives a natural
  comparison \(\Div^{\,n}(R) \to \Div^{\,n}_{\mathrm{Zar}}(R)\) compatible with the Picard
  class.

  This is the auxiliary carrier used by the statements in this section. The widened
  representability route uses the separate arbitrary-affine carrier described above. The reason
  for retaining this auxiliary carrier is
  Theorem~\ref{thm:isCertified_chart_obstruction}: the condition ``carries a global
  certificate'' is not local on \(\Spec R\), so \(\Div^{\,n}\) cannot be a Zariski sheaf on
  affine test algebras, whereas \(\Div^{\,n}_{\mathrm{Zar}}\) is local on the base by
  construction.
\end{definition}
\begin{proof}
  \leanok
  Divisor equality is an equivalence relation on all local-equation systems
  (\ref{lem:localEquations_divEq}) and preserves the property of being locally certified
  (\ref{def:isLocallyCertified}), so it restricts to an equivalence relation on the locally
  certified ones. The Picard class is constant on divisor-equality classes
  (\ref{lem:localEquations_divEq}), so it factors through the quotient. The comparison
  sends the class of a certified family \((d,A,c)\) to the class of \(d\), which is locally
  certified by the last part of \ref{def:isLocallyCertified}; it is well defined because
  both sides are quotients by the same relation, and it commutes with the Picard class
  because both are computed from \(d\).
\end{proof}

\begin{definition}[The widened arbitrary-affine divisor functor]
  \label{def:divFamZarAff}
  \provenancetag{CUSTOM}
  \lean{AlgebraicGeometry.CertifiedDivisorFamilyAff,
    AlgebraicGeometry.IsLocallyCertifiedAff, AlgebraicGeometry.DivFamZarAff,
    AlgebraicGeometry.DivFamZarAff.mk, AlgebraicGeometry.DivFamZarAff.mk_eq_mk_iff,
    AlgebraicGeometry.DivFamZarAff.picClass, AlgebraicGeometry.divFunctorAff,
    AlgebraicGeometry.divFunctorAff_obj, AlgebraicGeometry.divFunctorAff_map}
  \uses{def:relCurve, def:localEquations, lem:localEquations_divEq}
  \leanok
  For an affine test algebra \(R\), the R2 carrier is the quotient
  \[
    \Div^{\,n}_{\mathrm{Aff}}(R) \;:=\;
      \bigl\{\,d\text{ on }C_R\text{ with an arbitrary affine-open cover carrying a
      degree-}n\text{ certificate}\,\bigr\}\big/\!\sim,
  \]
  where the equivalence relation is equality of local divisors. The cover is part of the
  presentation and may vary with the divisor; no pair of pinned \(\PP^1\)-charts is required.
  Base change makes these values into the presheaf
  \(\operatorname{divFunctorAff}(C,n)\). Its representability target is therefore
  \[
    \operatorname{divFunctorAff}(C,n)\text{ is represented by }D,
  \]
  not a representation of the auxiliary \(\Div^{\,n}_{\mathrm{Zar}}\) carrier above.
  The carrier construction alone does not supply a representation witness.  The
  unconditional witnesses at degree \(g(C)\) and at the admissible parameter are recorded in
  Definitions~\ref{def:genusDivisorRepresentation} and
  \ref{def:admissibleDivisorRepresentation}.
\end{definition}

\section{The window embedding and the divisor scheme}
\label{sec:divscheme}

Grothendieck's route from a Quot-type functor to a scheme is: twist high enough that the
relevant \(H^1\) vanishes, push forward, and read the result as a point of a Grassmannian;
then cut out the image by the condition that the submodule be closed under multiplication
by sections of a further twist. The window engine in this section is deliberately the
auxiliary pinned-chart calculation for \(\Div^{\,n}_{\mathrm{Zar}}\); it is not an identification
with the widened carrier of Definition~\ref{def:divFamZarAff}.

Fix two exponents \(a\) with \(H^1(C_R, \OO(aF)) = 0\) universally --- concretely the two
\emph{pinned window exponents} \(M\) and \(M+s\) attached to \(\pi\) and the genus \(g\),
which are chosen past the bound at which the fibre twists have vanishing relative \(H^1\)
and the multiplication maps are surjective.

\begin{definition}[The affine Grassmannian functor]
  \label{def:grFunctorAff}
  \provenancetag{REFERENCE}
  \lean{AlgebraicGeometry.Grassmannian.grFunctorAff,
    AlgebraicGeometry.Grassmannian.grFunctorAffineEquiv}
  \leanok
  \source{nitsure-hilbert-quot}
  For a finite-dimensional \(k\)-vector space \(H\) and \(n \in \mathbb{N}\), the
  \emph{Grassmannian functor} \(\mathrm{Gr}(H, n)\) assigns to a \(k\)-algebra \(R\) the set
  of \(R\)-submodules \(K \subseteq R \otimes_k H\) whose quotient \((R \otimes_k H)/K\) is
  finite projective over \(R\) of constant rank \(n\) at every prime. A \(k\)-algebra map
  \(R \to R'\) acts by \(K \mapsto \operatorname{im}(R' \otimes_R K \to R' \otimes_k H)\).
  Corank rather than rank is the right index because it is what base-changes: a projective
  quotient of constant rank stays one.
\end{definition}

\begin{definition}[The window submodule and the embedding \(\varepsilon\)]
  \label{def:divisorWindow}
  \provenancetag{CUSTOM}
  \lean{AlgebraicGeometry.divisorWindow, AlgebraicGeometry.divisorWindow_eq_of_divEq,
    AlgebraicGeometry.DivFam.window, AlgebraicGeometry.divFamEps,
    AlgebraicGeometry.divFamEps_mk}
  \uses{def:certifiedDivisorFamily, def:divisorSections, def:fiberWeilDivisor,
    def:fiberTwist, lem:localEquations_divEq, def:grFunctorAff}
  \leanok
  \source{nitsure-hilbert-quot}
  Let \(a\) be an exponent at which the relative twist \(\OO(aF)\) on \(C_R\) has vanishing
  \(H^1\), so that base change identifies
  \[
    R \otimes_k H_a \;\xrightarrow{\ \sim\ }\; H^0\bigl(C_R, \OO(aF)\bigr).
  \]
  For a local-equation system \(d\) on \(C_R\), the \emph{window submodule}
  \[
    K_a(d) \;\subseteq\; R \otimes_k H_a
  \]
  is the preimage under this identification of the submodule of sections of \(\OO(aF)\)
  vanishing along \(d\) --- i.e.\ of \(H^0(C_R, \OO(aF - D))\), read inside the free
  window. It depends only on \(d\), not on any chart adaptation, and it is unchanged when
  \(d\) is replaced by a divisor-equal system; hence it descends to \(\Div^{\,n}(R)\) and to
  \(\Div^{\,n}_{\mathrm{Zar}}(R)\).

  Taking the two pinned exponents gives the \emph{window pair}
  \[
    \varepsilon(D) \;=\; \bigl(K_M(D),\, K_{M+s}(D)\bigr)
      \;\in\; \Bigl(R \otimes_k H_M \ \text{submodules}\Bigr)
        \times \Bigl(R \otimes_k H_{M+s}\ \text{submodules}\Bigr).
  \]
\end{definition}
\begin{proof}
  \leanok
  Vanishing along a divisor is a condition on the section, not on its presentation:
  a section \(\sigma\) of \(\OO(aF)\) lies in \(H^0(\OO(aF-D))\) iff on each member of the
  cover of \(d\) the local expression of \(\sigma\) is divisible by the equation of \(d\).
  If \(d \sim d'\), then on a common refinement the two equations differ by a unit at every
  point, and divisibility by an element is unchanged when the element is multiplied by a
  unit; divisibility is local on the cover, so the two vanishing submodules agree. Hence
  \(K_a(-)\) is constant on divisor-equality classes and factors through the quotients of
  \ref{def:certifiedDivisorFamily} and \ref{def:divFamZar}.
\end{proof}

\begin{lemma}[Equal rank at every stalk forces equality]
  \label{lem:rankAtStalk_equal_quotients}
  \provenancetag{CUSTOM}
  \lean{LinearMap.injective_of_surjective_of_rankAtStalk_eq,
    Submodule.eq_of_le_of_rankAtStalk_quotient_eq}
  \leanok
  Let \(R\) be a commutative ring.
  \begin{enumerate}
    \item A surjection \(\varphi \colon M \twoheadrightarrow N\) of \(R\)-modules with \(M\)
      finite projective and \(N\) projective, whose ranks at every prime agree, is an
      isomorphism.
    \item Consequently, if \(K_1 \subseteq K_2\) are submodules of \(M\) with \(M/K_1\)
      finite projective, \(M/K_2\) projective, and
      \(\operatorname{rank}_{\mathfrak p}(M/K_1) = \operatorname{rank}_{\mathfrak p}(M/K_2)\)
      for every prime \(\mathfrak p\), then \(K_1 = K_2\).
  \end{enumerate}
\end{lemma}
\begin{proof}
  \leanok
  (i) Since \(N\) is projective, \(\varphi\) splits: there is \(\ell \colon N \to M\) with
  \(\varphi \ell = \id_N\), and \(M \cong \ker\varphi \oplus N\). Hence \(\ker \varphi\) is
  a direct summand of the finite projective module \(M\), so it is itself finite
  projective, and rank is additive over direct sums:
  \(\operatorname{rank}_{\mathfrak p} \ker\varphi + \operatorname{rank}_{\mathfrak p} N =
  \operatorname{rank}_{\mathfrak p} M = \operatorname{rank}_{\mathfrak p} N\) for every
  \(\mathfrak p\), so \(\operatorname{rank}_{\mathfrak p}\ker\varphi = 0\) identically. A
  finite projective module of rank \(0\) at every prime is trivial (it is trivial after
  localisation at every prime, and triviality of a module is local). Hence
  \(\ker \varphi = 0\) and \(\varphi\) is an isomorphism.

  (ii) The canonical map \(M/K_1 \to M/K_2\) is surjective, and (i) applies to it, so it is
  injective. If \(x \in K_2\), its class in \(M/K_1\) maps to \(0\), hence is \(0\), i.e.\
  \(x \in K_1\). Thus \(K_2 \subseteq K_1\), and the reverse inclusion is the hypothesis.
\end{proof}

\begin{theorem}[The window identity]
  \label{thm:divisorWindow_eq_of_le}
  \provenancetag{CUSTOM}
  \lean{AlgebraicGeometry.divisorWindow_eq_of_le,
    AlgebraicGeometry.divisorWindow_eq_of_le_of_isCertified,
    AlgebraicGeometry.divFamEps_mk_eq_of_le}
  \uses{def:divisorWindow, def:isCertified, lem:rankAtStalk_equal_quotients,
    def:grFunctorAff, def:gluedColength}
  \leanok
  Let \(A\) be a finite chart adaptation of \(d\) over \(R\), certified of degree \(n\), let
  \(a\) be an exponent as in \ref{def:divisorWindow}, and assume the evaluation map
  \[
    R \otimes_k H_a \;\longrightarrow\; W(A)^{(a)}
  \]
  onto the \(\OO(aF)\)-twisted glued colength module is surjective. Then:
  \begin{enumerate}
    \item \((R \otimes_k H_a)/K_a(d) \cong W(A)^{(a)}\); in particular \(K_a(d)\) is a point
      of \(\mathrm{Gr}(H_a, n)\) over \(R\);
    \item any point \(x\) of \(\mathrm{Gr}(H_a, n)\) over \(R\) with \(x \subseteq K_a(d)\)
      equals \(K_a(d)\).
  \end{enumerate}
  Consequently, for a certified family whose two window submodules contain a given pair of
  Grassmannian points at the pinned exponents, \(\varepsilon\) of the family \emph{is} that
  pair.
\end{theorem}
\begin{proof}
  \leanok
  (i) By definition \(K_a(d)\) is the kernel of the evaluation map: a section of
  \(\OO(aF)\) vanishes along \(d\) exactly when its image in each chart colength module ---
  hence in their glued equalizer --- is zero. The evaluation is surjective by hypothesis,
  so the first isomorphism theorem gives the displayed isomorphism. The certificate makes
  \(W(A)^{(a)}\) finite projective of constant rank \(n\) at every prime: this is clause
  (c2) transported across the twist by \(\OO(aF)\), which is an invertible sheaf and hence
  does not change finiteness, projectivity or rank of the glued colength module. So the
  quotient of \(R \otimes_k H_a\) by \(K_a(d)\) is finite projective of constant rank
  \(n\), which is exactly membership in \(\mathrm{Gr}(H_a, n)(R)\)
  (\ref{def:grFunctorAff}).

  (ii) Both \((R\otimes_k H_a)/x\) and \((R \otimes_k H_a)/K_a(d)\) are finite projective of
  constant rank \(n\) --- the first because \(x\) is a Grassmannian point, the second by
  (i) --- so they have equal rank at every prime, and
  Lemma~\ref{lem:rankAtStalk_equal_quotients}(ii) applied to the inclusion
  \(x \subseteq K_a(d)\) gives \(x = K_a(d)\).

  The last sentence is (ii) applied at each of the two pinned exponents separately, the
  certificate slots at both exponents being supplied by the family's own certificate.
\end{proof}

\begin{theorem}[\(\varepsilon\) is injective on divisor classes]
  \label{thm:divFam_mono}
  \provenancetag{CUSTOM}
  \lean{AlgebraicGeometry.divFam_divEq_of_stalkIdeal_eq,
    AlgebraicGeometry.divFam_divEq_of_eps_eq_total}
  \uses{def:divisorWindow, def:certifiedDivisorFamily, lem:localEquations_divEq,
    def:isCertified}
  \leanok
  Assume the base-field normalisations \(h^0(\struct C) = 1\) and
  \(\chi(\struct C) = 1 - g\). Let \(G, G'\) be certified divisor families of degree \(g\)
  over an arbitrary \(k\)-algebra \(R\) --- no Noetherian hypothesis --- with
  \(\varepsilon(G) = \varepsilon(G')\). Then \(G\) and \(G'\) are divisor-equal; that is,
  \(\varepsilon\) is injective on \(\Div^{\,g}(R)\).
\end{theorem}
\begin{proof}
  \leanok
  \emph{Step 1: stalk ideals determine the class.} Two local-equation systems with the same
  ideal \((d_z)\struct{C_R,z}\) in every stalk are divisor-equal. Indeed, fix a point
  \(z\); the two equations generate the same stalk ideal, so each is a unit multiple of the
  other in the stalk, and a germ-level unit factor lifts to a unit on some neighbourhood of
  \(z\) contained in both members. Letting \(\mathcal W(z)\) be such a neighbourhood defines
  a pointed cover refining both, on which the equations differ by a unit at every point ---
  which is divisor equality (\ref{lem:localEquations_divEq}). Descending through the
  quotient of \ref{def:certifiedDivisorFamily} gives equality in \(\Div^{\,g}(R)\).

  \emph{Step 2: the window generates.} For a certified family the first window generates the
  divisor: at every point \(z\) of \(C_R\) the stalk ideal of \(d\) is generated by the
  germs at \(z\) of the sections in \(K_M(d)\), divided by a local generator of
  \(\OO(MF)\). This is where the exponent \(M\) is used: it is past the bound at which
  \(\OO(MF - D)\) is relatively globally generated for every degree-\(g\) certified family,
  and global generation of \(\OO(MF-D)\) says exactly that the image of
  \(H^0(\OO(MF-D)) \otimes \struct{} \to \OO(MF)\) is the subsheaf \(\OO(MF-D)\), whose
  local generator is (a local generator of \(\OO(MF)\) times) the equation of \(D\). The
  normalisations \(h^0 = 1\), \(\chi = 1-g\) are what pin the numerical bound.

  \emph{Conclusion.} If \(\varepsilon(G) = \varepsilon(G')\) then in particular
  \(K_M(G) = K_M(G')\), so by Step 2 the two systems have equal stalk ideals at every
  point, and by Step 1 they are divisor-equal.
\end{proof}

\begin{definition}[The divisor scheme \(\mathbf{Div}\)]
  \label{def:divScheme}
  \provenancetag{ADAPTED}
  \lean{AlgebraicGeometry.DivScheme, AlgebraicGeometry.divSchemeι,
    AlgebraicGeometry.divSchemeOver, AlgebraicGeometry.divScheme_hom_ext,
    AlgebraicGeometry.divCarveMul, AlgebraicGeometry.divCarveIdeal}
  \uses{def:grFunctorAff, def:divisorSections, def:fiberWeilDivisor}
  \leanok
  \source{nitsure-hilbert-quot, kleiman-picard}
  Choose \(k\)-bases of \(H_M\) and \(H_{M+s}\), giving the Grassmannian functors
  \(\mathrm{Gr}(H_M, g)\) and \(\mathrm{Gr}(H_{M+s}, g)\) their standard affine charts
  indexed by coordinate subsets, and let \(\mathrm{Gr}\text{-pair}\) be the product of the
  two chart-glued Grassmannian schemes. For \(\sigma \in H_s = H^0(C, \OO(sF))\), multiplication
  by \(\sigma\) is a \(k\)-linear map \(\mu_\sigma \colon H_M \to H_{M+s}\).

  The \emph{divisor scheme} \(\mathbf{Div}\) is the closed subscheme of
  \(\mathrm{Gr}\text{-pair}\) cut out, chart by chart, by the ideal generated by the matrix
  entries of the composites
  \[
    K_1 \;\hookrightarrow\; R \otimes_k H_M \;\xrightarrow{\ \id \otimes \mu_\sigma\ }\;
      R \otimes_k H_{M+s} \;\twoheadrightarrow\; (R \otimes_k H_{M+s})/K_2 ,
    \qquad \sigma \in H_s ,
  \]
  where \((K_1, K_2)\) is the tautological pair on the chart. Its \(R\)-points are therefore
  the pairs \((K_1, K_2)\) of Grassmannian points of corank \(g\) satisfying
  \(\mu_\sigma(K_1) \subseteq K_2\) for every \(\sigma \in H_s\). It comes with a closed
  immersion \(\mathbf{Div} \hookrightarrow \mathrm{Gr}\text{-pair}\); in particular
  morphisms into \(\mathbf{Div}\) are determined by their composites with that immersion.
  Placed over \(\Spec k\) by its structure morphism, it is an object
  \(\mathbf{Div}_{/k}\) of \(\Over(\Spec k)\).
\end{definition}
\begin{proof}
  \leanok
  The description of the \(R\)-points is the standard reading of a determinantal-type carve
  on a Grassmannian chart: on the chart the tautological submodules are free with an
  explicit basis, the composite above is an \(R\)-linear map between free modules, and a map
  between free modules vanishes exactly when all its matrix entries do. Hence a \(k\)-algebra
  map from the chart ring kills the carve ideal iff the base-changed composite is zero for
  every \(\sigma\), i.e.\ iff \(\mu_\sigma(K_1) \subseteq K_2\). The carve ideals on
  overlapping charts agree because the condition they express is chart-independent, so they
  glue to a quasi-coherent ideal sheaf and the closed subschemes glue. Determination of
  morphisms by their composite with the immersion holds because a closed immersion is a
  monomorphism.
\end{proof}

The relation of \ref{def:divScheme} to \ref{def:divisorWindow} is the reason for the carve:
if \(D\) is an effective divisor then \(\sigma \cdot H^0(\OO(MF - D)) \subseteq
H^0(\OO((M+s)F - D))\) for every \(\sigma \in H^0(\OO(sF))\), so \(\varepsilon(D)\) is an
\(R\)-point of \(\mathbf{Div}\). The content of the construction is the converse direction,
that a point of \(\mathbf{Div}\) determines a divisor; that is the forward map of
Section~\ref{sec:divrep}, which is not built.

\section{The classifying morphism}
\label{sec:divclassify}

This section is the backward half of representability: a divisor class over an affine test
determines a morphism to \(\mathbf{Div}\), uniquely, and the class is recoverable from the
morphism. Both directions are proved.

\begin{definition}[The characterizing clause of a classifying morphism]
  \label{def:isDivRepClassify}
  \provenancetag{CUSTOM}
  \lean{AlgebraicGeometry.IsDivRepClassify}
  \uses{def:divFamZar, def:divScheme, def:divisorWindow, def:grFunctorAff}
  \leanok
  Let \(S\) be a \(k\)-algebra, \(D \in \Div^{\,g}_{\mathrm{Zar}}(S)\), and \(v \colon \Spec
  S \to \mathbf{Div}\) a morphism. Say \(v\) \emph{classifies} \(D\) if for every
  \(k\)-algebra \(T\) under \(S\), every certified family \(G\) over \(T\) whose class is
  the restriction of \(D\) to \(T\), and every presentation of \(\varepsilon(G)\) as a point
  of a chart of \(\mathrm{Gr}\text{-pair}\) --- that is, a \(k\)-algebra map \(w\) from the
  chart ring to \(T\) carrying the tautological pair to \(\varepsilon(G)\) --- the square
  \[
    \Spec T \longrightarrow \Spec S \xrightarrow{\ v\ } \mathbf{Div}
      \hookrightarrow \mathrm{Gr}\text{-pair}
  \]
  agrees with the chart morphism \(\Spec T \xrightarrow{w} \mathrm{Gr}\text{-pair}\).

  Quantifying over all tests \(T\), all certified representatives and all chart
  presentations makes the clause stable under refinement of covers and change of
  representative by construction.
\end{definition}

\begin{theorem}[Existence and uniqueness of the classifying morphism]
  \label{thm:divClassifyZar}
  \provenancetag{CUSTOM}
  \lean{AlgebraicGeometry.exists_isDivRepClassify,
    AlgebraicGeometry.isDivRepClassify_unique, AlgebraicGeometry.divClassifyZar,
    AlgebraicGeometry.divRepClassifyZar,
    AlgebraicGeometry.divRepClassifyZar_isDivRepClassify,
    AlgebraicGeometry.divRepClassifyZar_eq_of_isDivRepClassify}
  \uses{def:isDivRepClassify, def:divFamZar, def:divScheme, thm:divisorWindow_eq_of_le,
    def:isLocallyCertified}
  \leanok
  For every \(k\)-algebra \(S\) and every \(D \in \Div^{\,g}_{\mathrm{Zar}}(S)\) there is a
  unique morphism \(\Spec S \to \mathbf{Div}\) classifying \(D\) in the sense of
  \ref{def:isDivRepClassify}. It lifts canonically to a morphism
  \[
    \gamma_S(D) \colon \Spec S \longrightarrow \mathbf{Div}_{/k}
  \]
  in \(\Over(\Spec k)\), and any morphism of \(\Over(\Spec k)\) whose underlying scheme
  morphism classifies \(D\) equals \(\gamma_S(D)\).
\end{theorem}
\begin{proof}
  \leanok
  \emph{Existence.} By \ref{def:isLocallyCertified} there is a finite family
  \(r_1,\dots,r_m \in S\) generating the unit ideal such that over each \(S_{r_p}\) the
  restriction of \(D\) has a certified representative \(G_p\); refining that family further
  we may also assume that \(\varepsilon(G_p)\) admits a chart presentation \(w_p\) over
  \(S_{r_p}\), since the Grassmannian charts cover \(\mathrm{Gr}\text{-pair}\) and a point
  of a scheme covered by opens lies in one of them after a further localisation of the
  base. Over each \(S_{r_p}\) the chart presentation kills the carve ideal --- by the
  multiplication inclusion recorded after \ref{def:divScheme} --- so it factors through the
  closed immersion of \ref{def:divScheme}, giving \(v_p \colon \Spec S_{r_p} \to
  \mathbf{Div}\) satisfying the clause for \(G_p\).

  The \(v_p\) agree on overlaps: on \(\Spec S_{r_p r_q}\) both restrictions satisfy the
  clause for the restriction of \(D\), and the uniqueness argument below (applied over the
  ring \(S_{r_p r_q}\)) identifies them. Gluing along the affine open cover \(\{\Spec
  S_{r_p}\}\) of \(\Spec S\) yields \(v\), and the clause for \(v\) follows from the clause
  for each \(v_p\) because it may be tested after pulling back to the members of a cover.

  \emph{Uniqueness.} Let \(v, v'\) both classify \(D\). Take a family \(r_p\) and data
  \(G_p, w_p\) as above. Restricting to \(\Spec S_{r_p}\), the clause applied to \(v\) and
  to \(v'\) with the same \((T, G, w) = (S_{r_p}, G_p, w_p)\) gives that both composites
  \(\Spec S_{r_p} \to \mathbf{Div} \hookrightarrow \mathrm{Gr}\text{-pair}\) equal the same
  chart morphism. Since \(\mathbf{Div} \hookrightarrow \mathrm{Gr}\text{-pair}\) is a
  monomorphism (\ref{def:divScheme}), the two restrictions to \(\Spec S_{r_p}\) agree; as
  the \(\Spec S_{r_p}\) cover \(\Spec S\), \(v = v'\).

  \emph{The lift.} The composite of \(v\) with the immersion is, over each member of the
  cover, a chart morphism, and every chart of \(\mathrm{Gr}\text{-pair}\) lies over
  \(\Spec k\) compatibly; so \(v\) commutes with the structure morphisms to \(\Spec k\) and
  underlies a morphism \(\gamma_S(D)\) of \(\Over(\Spec k)\). A morphism of
  \(\Over(\Spec k)\) is determined by its underlying scheme morphism, so the last assertion
  is uniqueness again.
\end{proof}

\begin{theorem}[Separation of divisor classes]
  \label{thm:eq_of_isDivRepClassify}
  \provenancetag{CUSTOM}
  \lean{AlgebraicGeometry.eq_of_isDivRepClassify,
    AlgebraicGeometry.divRepClassifyZar_injective}
  \uses{thm:divClassifyZar, def:isDivRepClassify, def:divFamZar, thm:divFam_mono,
    def:isLocallyCertified}
  \leanok
  If \(D_0, D_1 \in \Div^{\,g}_{\mathrm{Zar}}(S)\) are classified by the \emph{same}
  morphism \(v \colon \Spec S \to \mathbf{Div}\), then \(D_0 = D_1\). Equivalently,
  \(D \mapsto \gamma_S(D)\) is injective.
\end{theorem}
\begin{proof}
  \leanok
  Choose, as in \ref{thm:divClassifyZar}, unit-ideal families \(c^{(0)}\) and \(c^{(1)}\) in
  \(S\) over whose localisations \(D_0\), respectively \(D_1\), have certified
  representatives with chart-presented \(\varepsilon\)-pairs. The products
  \(c^{(0)}_p c^{(1)}_q\) again generate the unit ideal (a product of two unit ideals is the
  unit ideal), so it suffices to prove that \(D_0\) and \(D_1\) have the same restriction to
  each \(S_{c^{(0)}_p c^{(1)}_q}\); membership in \(\Div^{\,g}_{\mathrm{Zar}}\) is a
  quotient of local-equation systems and equality of classes is local on \(\Spec S\).

  Fix such a common localisation \(B\). Over \(B\) we have certified representatives
  \(G_0\) of \(D_0|_B\) and \(G_1\) of \(D_1|_B\), together with chart presentations \(w_0\)
  of \(\varepsilon(G_0)\) and \(w_1\) of \(\varepsilon(G_1)\). Applying the clause of
  \ref{def:isDivRepClassify} to the single morphism \(v\), once with \((G_0, w_0)\) and once
  with \((G_1, w_1)\), shows that \(w_0\) and \(w_1\) present the \emph{same} morphism
  \(\Spec B \to \mathrm{Gr}\text{-pair}\); hence the two tautological images agree, i.e.
  \[
    \varepsilon(G_0) = \varepsilon(G_1)
  \]
  as pairs of submodules of the free windows over \(B\). By the mono theorem
  (\ref{thm:divFam_mono}) two certified families of degree \(g\) over \(B\) with equal
  \(\varepsilon\)-pairs are divisor-equal, so \(G_0\) and \(G_1\) define the same class;
  hence \(D_0|_B = D_1|_B\).

  Note that this argument consumes neither a certificate over the chart ring of
  \(\mathbf{Div}\) nor any universal family: the only certified objects are the local
  representatives supplied by \ref{def:isLocallyCertified}.
\end{proof}

\section{Affine representability, and what remains}
\label{sec:divrep}

Everything above is the backward direction: from divisor classes to morphisms. Turning
\(\mathbf{Div}\) into a representing object requires the forward direction --- a universal
divisor family on \(\mathbf{Div}\), from which a divisor class can be pulled back along any
test morphism. This section states exactly what that datum is, shows that it suffices, and
records that it has not been constructed.

\begin{definition}[The affine pullback package]
  \label{def:divRepAffinePullback}
  \provenancetag{CUSTOM}
  \lean{AlgebraicGeometry.divRepPullAt, AlgebraicGeometry.divRepPullAt_comp,
    AlgebraicGeometry.DivRepChartFamily.IsCompatible,
    AlgebraicGeometry.DivRepAffinePullback}
  \uses{def:divFamZar, def:divScheme, def:isDivRepClassify}
  \leanok
  An \emph{affine pullback package} for \(\mathbf{Div}\) consists of, for every
  \(k\)-algebra \(S\), a map
  \[
    \rho_S \colon \Hom_{\Over(\Spec k)}\bigl(\Spec S, \mathbf{Div}_{/k}\bigr)
      \longrightarrow \Div^{\,g}_{\mathrm{Zar}}(S),
  \]
  subject to two conditions: \(\rho_S(v)\) is classified by \(v\) in the sense of
  \ref{def:isDivRepClassify}, for every \(v\); and \(\rho\) is natural, \(\rho_B(\varphi
  \circ v) = \varphi_* \rho_A(v)\) for every \(k\)-algebra map \(\varphi \colon A \to B\).

  The canonical way to produce such a \(\rho\) is from a \emph{chart family}: a certified
  divisor family \(U_{ij}\) over the coordinate ring of each affine chart of
  \(\mathbf{Div}\), pulled back along a chart presentation of the given \(v\); this is well
  defined precisely when the family is \emph{compatible}, meaning that two chart
  presentations inducing the same morphism to \(\mathbf{Div}\) pull \(U\) back to the same
  class in \(\Div^{\,g}_{\mathrm{Zar}}\).
\end{definition}

\begin{theorem}[The affine equivalence]
  \label{thm:divRepAffinePullback_equiv}
  \provenancetag{CUSTOM}
  \lean{AlgebraicGeometry.pull_classify_of_isDivRepClassify_pull,
    AlgebraicGeometry.DivRepAffinePullback.ofPull,
    AlgebraicGeometry.DivRepAffinePullback.equiv,
    AlgebraicGeometry.DivRepAffinePullback.equiv_naturality}
  \uses{def:divRepAffinePullback, thm:divClassifyZar, thm:eq_of_isDivRepClassify}
  \leanok
  An affine pullback package \(\rho\) makes \(\mathbf{Div}_{/k}\) represent
  \(\Div^{\,g}_{\mathrm{Zar}}\) on affine tests: for every \(k\)-algebra \(S\) the maps
  \(\rho_S\) and \(\gamma_S\) (\ref{thm:divClassifyZar}) are mutually inverse bijections
  \[
    \Hom_{\Over(\Spec k)}\bigl(\Spec S, \mathbf{Div}_{/k}\bigr)
      \;\simeq\; \Div^{\,g}_{\mathrm{Zar}}(S),
  \]
  natural in \(S\).
\end{theorem}
\begin{proof}
  \leanok
  \(\gamma_S \circ \rho_S = \id\): given \(v\), the class \(\rho_S(v)\) is classified by
  \(v\) by hypothesis, and by the uniqueness clause of \ref{thm:divClassifyZar} any
  morphism classifying a class \emph{is} the classifying morphism of that class; hence
  \(\gamma_S(\rho_S(v)) = v\).

  \(\rho_S \circ \gamma_S = \id\): given \(D\), both \(\rho_S(\gamma_S(D))\) and \(D\) are
  classified by the single morphism \(\gamma_S(D)\) --- the first by the hypothesis on
  \(\rho\), the second by the characterizing property of \(\gamma_S\)
  (\ref{thm:divClassifyZar}). By separation (\ref{thm:eq_of_isDivRepClassify}) two classes
  classified by the same morphism are equal, so \(\rho_S(\gamma_S(D)) = D\). (In particular
  this round-trip law need not be assumed in \ref{def:divRepAffinePullback}: it is a
  consequence of the clause law and separation.)

  Naturality of the bijection in \(S\) is the naturality assumed of \(\rho\).
\end{proof}

\begin{theorem}[The auxiliary fixed-window universal family --- unproved]
  \label{thm:divRepUniversalFamily}
  \provenancetag{CUSTOM}
  \uses{def:divRepAffinePullback, def:divScheme, def:certifiedDivisorFamily,
    thm:divisorWindow_eq_of_le, thm:isCertified_chart_obstruction}
  \notready
  There exists an affine pullback package for \(\mathbf{Div}\) in the sense of
  \ref{def:divRepAffinePullback}; equivalently, there is a compatible chart family --- a
  certified divisor family \(U_{ij}\) of degree \(g\) over the coordinate ring of each
  affine chart of \(\mathbf{Div}\), whose window pair is the tautological pair of that
  chart, and whose pullbacks along two presentations of one morphism agree in
  \(\Div^{\,g}_{\mathrm{Zar}}\).

  This statement is not proved. It is the forward direction of representability: given a
  point \((K_1, K_2)\) of \(\mathbf{Div}\) over \(R\), one must produce the divisor whose
  windows are \(K_1\) and \(K_2\). The intended construction takes a generating family of
  \(K_1\) inside \(R \otimes_k H_M\) and reads its common vanishing locus as an effective
  divisor, with local equations extracted chart by chart; the two obstacles are that this
  produces the local equations only after a Zariski localisation of \(R\) --- which is why
  the target is \(\Div^{\,g}_{\mathrm{Zar}}\) and not \(\Div^{\,g}\), and why the global
  certificate of \ref{def:isCertified} is unavailable
  (\ref{thm:isCertified_chart_obstruction}) --- and that the resulting family's second
  window must be shown to be exactly \(K_2\), for which the containment
  \(K_2 \subseteq K_{M+s}\) has to be upgraded from the fibrewise statement to a relative
  one, an upgrade that fails over nonreduced base rings if attempted fibrewise.

  Everything in Sections~\ref{sec:divfam}--\ref{sec:divclassify} is independent of this
  statement. This is a conditional result for the pinned-window auxiliary carrier, not the
  current R2 representability target \(\operatorname{divFunctorAff}(C,n)\).
\end{theorem}

\section{The rank-one Abel chart and descent}
\label{sec:rankOneAbelDescent}

The Abel fibre over a line-bundle class is its complete linear system.  At degree \(g(C)\),
the locus with \(H^1=0\) has one-dimensional space of sections, so the fibre is a point and
the Abel map can be inverted without choosing a generator.  At a larger degree the same
fibre generally has positive dimension.  Thus the represented degree-\(g(C)\) divisor
functor supplies the chart, while the arbitrary-degree construction remains auxiliary.

\subsection{The represented sources and the fibre obstruction}

\begin{definition}[The represented genus-degree divisor source]
  \label{def:genusDivisorRepresentation}
  \provenancetag{CUSTOM}
  \lean{AlgebraicGeometry.divRepAffGenusScheme,
    AlgebraicGeometry.divFunctorAff_genus_representableBy}
  \uses{def:divFamZarAff}
  \leanok
  There is a scheme \(D_g\) over \(\Spec k\) representing
  \(\operatorname{divFunctorAff}(C,g(C))\).  In particular, a morphism
  \(T\to D_g\) is naturally the same as a relative effective divisor of degree \(g(C)\) on
  \(C_T/T\), presented on an arbitrary affine-open cover.
\end{definition}

\begin{definition}[The represented admissible divisor source]
  \label{def:admissibleDivisorRepresentation}
  \provenancetag{CUSTOM}
  \lean{AlgebraicGeometry.divRepAffAdmissibleParameter,
    AlgebraicGeometry.divRepAffAdmissibleScheme,
    AlgebraicGeometry.divFunctorAff_admissible_representableBy}
  \uses{def:divFamZarAff}
  \leanok
  There is a parameter \(n_{\mathrm{adm}}\), determined by the curve data, and a scheme
  \(D_{\mathrm{adm}}\) over \(\Spec k\) representing
  \(\operatorname{divFunctorAff}(C,n_{\mathrm{adm}})\), without a rational-point or further
  representability hypothesis.
\end{definition}

\begin{definition}[The represented genus-degree Abel map]
  \label{def:genusAbelMap}
  \provenancetag{CUSTOM}
  \lean{AlgebraicGeometry.abelDivAffTrans, AlgebraicGeometry.abelDivAffSigma,
    AlgebraicGeometry.abelDivAffGenusSigma}
  \uses{def:genusDivisorRepresentation, def:picDegLayerFunctor}
  \leanok
  Sending a relative effective divisor to the class of its associated invertible sheaf gives
  a natural transformation
  \[
    a_g \colon h_{D_g} \longrightarrow \mathrm{pic}^{g(C)}(C).
  \]
  More generally the same construction exists in every degree and is compatible with the
  Sigma-extension from the slice over \(\Spec k\) to the big site.
\end{definition}

\begin{theorem}[The higher-degree Abel map is not an open chart]
  \label{thm:highDegreeAbel_notOpenImmersion}
  \provenancetag{ADAPTED}
  \alignmentcheck{TO CHECK: Milne page~0109 is under the standing hypothesis \(g(C)>0\), while
  this higher-degree Abel statement is packaged for all genera; verify the genus-zero scope.}
  \lean{AlgebraicGeometry.not_isOpenImmersion_abelSigmaChartAff_of_genus_lt_degree}
  \uses{def:divFamZarAff, thm:riemann_inequality_curve}
  \leanok
  Let \(n>g(C)\), fix a representation of the degree-\(n\) divisor functor, and fix any
  degree-calibrated translation of its Abel map.  If that divisor functor has a point after a
  field extension, then the corresponding unrestricted Abel map is not an open immersion of
  presheaves.
  \source{abelian-varieties:page-0109}
\end{theorem}
\begin{proof}
  Let \(D\) be such a divisor.  Riemann--Roch gives
  \(h^0(\mathcal O(D))\ge n+1-g(C)\ge2\).  Two independent sections determine two distinct
  effective divisors in the complete linear system of \(D\): if their divisors were equal,
  their ratio would be a global unit on the proper geometrically integral curve and hence a
  scalar, contrary to independence.  The two divisors therefore have the same Abel class but
  are unequal.  Thus the Abel map is not a monomorphism, whereas every open immersion is a
  monomorphism.
\end{proof}

\subsubsection{Admissible-degree Abel properties}

The admissible-degree Abel map remains useful for divisor geometry and for detecting the
relative-linear-equivalence relation.  The properties below do not by themselves assert the
existence of a scheme quotient.

\begin{definition}[The concrete admissible Abel map]
  \label{def:admissibleAbelMap}
  \provenancetag{CUSTOM}
  \lean{AlgebraicGeometry.admissibleAbelChartIndex,
    AlgebraicGeometry.admissibleAbelChartDivisor,
    AlgebraicGeometry.admissibleAbelTrans,
    AlgebraicGeometry.abelSigmaChartAffAdmissible_eq_sigmaExtension_admissibleAbelTrans}
  \uses{def:admissibleDivisorRepresentation, def:pic0SigmaFunctor}
  \leanok
  Choosing the admissible chart index and reference divisor internally gives a natural
  transformation from the represented admissible divisor functor to
  \(\mathrm{pic}^0(C)\).  Its Sigma-extension is the corresponding map on the big site.
\end{definition}

\begin{theorem}[The admissible source is project-local]
  \label{thm:admissibleDivisorSource_projectLocal}
  \provenancetag{ADAPTED}
  \lean{AlgebraicGeometry.divRepAffAdmissibleEmbedding,
    AlgebraicGeometry.isClosedImmersion_divRepAffAdmissibleEmbedding,
    AlgebraicGeometry.divRepAffAdmissibleGrassmannianCover,
    AlgebraicGeometry.quasiCompact_divRepAffAdmissibleScheme,
    AlgebraicGeometry.locallyOfFiniteType_divRepAffAdmissibleScheme,
    AlgebraicGeometry.isSeparated_divRepAffAdmissibleScheme}
  \uses{def:admissibleDivisorRepresentation, def:divScheme}
  \leanok
  The source \(D_{\mathrm{adm}}\) has a closed immersion into a product of two
  Grassmannians.  Consequently it is quasi-compact, locally of finite type, and separated,
  and it has a finite affine cover obtained by pulling back standard affine Grassmannian
  charts.
  \source{kleiman-picard}
\end{theorem}
\begin{proof}
  The two universal subbundles associated to an admissible divisor give the classifying map
  to the product of Grassmannians.  Equality of the classified subbundles recovers the
  divisor datum, so the map is a closed immersion.  Each Grassmannian is projective and has a
  finite standard affine cover.  A closed subscheme of their product is separated, locally
  of finite type, and quasi-compact, and the inverse images of the finitely many product
  charts give the asserted affine cover.
\end{proof}

\begin{theorem}[Kernel of the admissible Abel map]
  \label{thm:admissibleAbel_kernel}
  \provenancetag{ADAPTED}
  \lean{AlgebraicGeometry.admissibleAbelTrans_app_eq_iff_forall_picClass_div_mem_picFromBase}
  \uses{def:admissibleAbelMap, def:relPic}
  \leanok
  Two admissible divisor families have the same Abel value exactly when, on every affine open
  of the test scheme, the ratio of their Picard classes is pulled back from that affine open.
  Thus its kernel is relative linear equivalence, including the quotient by line bundles from
  the base.
  \source{kleiman-picard}
\end{theorem}
\begin{proof}
  The two Abel values are obtained from the divisor classes by multiplying by the same fixed
  reference class, so equality is unchanged after cancelling that class.  Equality in the
  relative Picard functor means precisely that, locally on the test, the two representing
  invertible sheaves differ by the pullback of an invertible sheaf on the test.  Restricting
  to an affine-open cover yields the stated condition.  Conversely those local pullback
  classes identify the two relative classes and hence the two Abel values.
\end{proof}

\begin{theorem}[Fibre dimension at the admissible degree]
  \label{thm:admissibleAbel_fibreRank}
  \provenancetag{ADAPTED}
  \alignmentcheck{TO CHECK: Milne page~0109 is under the standing hypothesis \(g(C)>0\), while
  this fibre-rank statement is packaged for all genera; verify the genus-zero scope.}
  \lean{AlgebraicGeometry.subsingleton_h1_of_divRepAffAdmissibleParameter_le_deg,
    AlgebraicGeometry.h0_divisorSheaf_eq_divRepAffAdmissibleParameter_add_one_sub_genus,
    AlgebraicGeometry.h0_sub_one_eq_divRepAffAdmissibleParameter_sub_genus}
  \uses{def:admissibleDivisorRepresentation, thm:h0_eq_deg_add_chi,
    thm:chi_moduleKSheaf_curve}
  \leanok
  On every field fibre, a divisor of exact admissible degree has vanishing \(H^1\), section
  rank \(n_{\mathrm{adm}}+1-g(C)\), and complete linear system of dimension
  \(n_{\mathrm{adm}}-g(C)\).
  \source{abelian-varieties:page-0109}
\end{theorem}
\begin{proof}
  The admissible bound gives \(H^1(\mathcal O(D))=0\).  Riemann--Roch therefore gives
  \(h^0(\mathcal O(D))=\deg D+1-g(C)=n_{\mathrm{adm}}+1-g(C)\).  The complete linear
  system is the projective space of nonzero sections modulo scalars, so its dimension is one
  less, namely \(n_{\mathrm{adm}}-g(C)\).
\end{proof}

\begin{definition}[The admissible Abel map after big-etale sheafification]
  \label{def:admissibleAbelEtaleSheafMap}
  \provenancetag{CUSTOM}
  \lean{AlgebraicGeometry.Scheme.subcanonical_etaleTopology,
    AlgebraicGeometry.pic0SigmaEtaleSheafification,
    AlgebraicGeometry.admissibleAbelEtaleSourceIso,
    AlgebraicGeometry.admissibleAbelEtaleSheafMap}
  \uses{def:admissibleDivisorRepresentation, def:admissibleAbelMap}
  \leanok
  Subcanonicity identifies the sheafification of the represented source with the big-etale
  Yoneda sheaf of \(D_{\mathrm{adm}}\).  Sheafifying the admissible Abel transformation gives
  a morphism from this Yoneda sheaf to the big-etale sheafification of the Sigma Picard
  presheaf.
\end{definition}

\begin{theorem}[Kernel after big-etale sheafification]
  \label{thm:admissibleAbel_etaleKernel}
  \provenancetag{CUSTOM}
  \lean{AlgebraicGeometry.admissibleAbelEtaleSheafMap_app_eq_iff_equalizerSieve_mem,
    AlgebraicGeometry.admissibleAbelEtaleRawKernelSieve_apply_iff_forall_picClass_div_mem_picFromBase}
  \uses{thm:admissibleAbel_kernel, def:admissibleAbelEtaleSheafMap}
  \leanok
  Two sections have equal images under the sheafified Abel map exactly when the sieve on which
  their pulled-back divisor classes differ by a line bundle from the base is covering.
\end{theorem}
\begin{proof}
  Equality in a sheafification holds precisely when the equalizer sieve of the two raw
  sections is covering.  By Theorem~\ref{thm:admissibleAbel_kernel}, an arrow belongs to that
  equalizer exactly when, on every affine open after pullback, the ratio of the two divisor
  classes comes from the base.  Substituting this description into the equalizer criterion
  gives the assertion.
\end{proof}

\begin{theorem}[Etale-local surjectivity of the admissible Abel map]
  \label{thm:admissibleAbel_etaleLocallySurjective}
  \provenancetag{ADAPTED}
  \lean{AlgebraicGeometry.isLocallySurjective_abelSigmaChartAffAdmissible}
  \uses{def:admissibleAbelMap, def:admissibleAbelEtaleSheafMap,
    thm:admissibleAbel_fibreRank}
  \leanok
  Every degree-zero Picard class lifts to an admissible divisor family after an etale cover of
  the test scheme.
  \source{kleiman-picard}
\end{theorem}
\begin{proof}
  After an etale cover, represent the Picard class by an invertible sheaf and apply the fixed
  admissible translation, obtaining an invertible sheaf of degree \(n_{\mathrm{adm}}\).
  Its first cohomology vanishes by Theorem~\ref{thm:admissibleAbel_fibreRank}, so its
  pushforward is locally free, commutes with base change, and has positive rank.  The
  projective bundle of rank-one quotients of the dual pushforward is smooth and surjective
  over the test.  It therefore has a section after a further etale cover.  The corresponding
  nowhere-zero fibrewise section of the invertible sheaf cuts out a relative effective
  divisor whose translated Abel class is the original class.
\end{proof}

\begin{theorem}[The sheafified Abel map is the coequalizer of its kernel pair]
  \label{thm:admissibleAbel_etaleCoequalizer}
  \provenancetag{CUSTOM}
  \lean{AlgebraicGeometry.admissibleAbelEtaleSheafMap_isLocallySurjective,
    AlgebraicGeometry.admissibleAbelEtaleTargetCoequalizer}
  \uses{def:admissibleAbelEtaleSheafMap,
    thm:admissibleAbel_etaleLocallySurjective}
  \leanok
  The sheafified admissible Abel map is the coequalizer of its pullback kernel pair in the
  category of big-etale sheaves.
\end{theorem}
\begin{proof}
  Etale-local surjectivity says exactly that the sheafified map is an epimorphism.  A category
  of sheaves of sets is a topos, and every epimorphism in a topos is effective: it is the
  coequalizer of the two projections from its pullback kernel pair.  Apply this fact to the
  sheafified Abel map.
\end{proof}

\begin{theorem}[Geometry of the genus-degree divisor source]
  \label{thm:genusDivisorSource_geometry}
  \provenancetag{ADAPTED}
  \uses{def:genusDivisorRepresentation}
  \notready
  The scheme \(D_g\) is projective and geometrically irreducible over \(k\).
  \source{kleiman-picard, abelian-varieties:page-0100,
    abelian-varieties:page-0107}
\end{theorem}
\begin{proof}
  The degree-\(g(C)\) relative effective-divisor functor is represented by the symmetric power
  \(C^{(g(C))}\): a tuple of points is sent to the sum of its graph divisors, and the universal
  property of the symmetric quotient identifies this with every relative effective divisor
  after base change.  The quotient map \(C^{g(C)}\to C^{(g(C))}\) is finite and surjective.
  Since \(C^{g(C)}\) is projective and geometrically irreducible, its finite symmetric quotient
  is projective, and its underlying space remains geometrically irreducible.  Uniqueness of a
  representing object identifies this symmetric power with \(D_g\).
\end{proof}

The admissible representer in Definition~\ref{def:admissibleDivisorRepresentation} provides
arbitrary-degree divisor geometry, but
Theorem~\ref{thm:highDegreeAbel_notOpenImmersion} precludes using its raw Abel map as the
chart below.  A quotient of that map would require, as additional mathematics, a represented
kernel relation, an effective scheme quotient compatible with base change, and a proof that
its Yoneda sheaf is the Picard sheaf.  Those assertions do not follow from the admissible-degree
properties above.

\subsection{The intrinsic rank-one locus}

\begin{definition}[A rank-one local presentation of a Picard class]
  \label{def:picRankOneLocalPresentation}
  \provenancetag{CUSTOM}
  \lean{AlgebraicGeometry.PicRankOneLocalPresentation}
  \uses{def:picDegLayerFunctor}
  \leanok
  Let \(T\) be a \(k\)-scheme, let \(f_T\colon C_T\to T\), and let
  \(\lambda\in\mathrm{pic}^{g(C)}(C,T)\).  A \emph{rank-one local presentation of
  \(\lambda\)} consists of an etale cover \(p_i\colon T_i\to T\) and invertible sheaves
  \(\mathcal L_i\) on \(C_{T_i}\) whose classes are exactly
  \(p_i^*\lambda\), such that
  \[
    R^1f_{i*}\mathcal L_i=0,
    \qquad \mathcal N_i:=f_{i*}\mathcal L_i\text{ is invertible},
  \]
  and the base-change map
  \(q^*\mathcal N_i\to f'_*q_C^*\mathcal L_i\) is an isomorphism for every
  \(q\colon T'\to T_i\).  The equation
  \([\mathcal L_i]=p_i^*\lambda\) is part of the data; a line bundle unrelated to
  \(\lambda\) is not a presentation of it.  The section used below is fixed to be the
  counit \(f_i^*\mathcal N_i\to\mathcal L_i\); no local generator of \(\mathcal N_i\) is
  included in the presentation.
\end{definition}

\begin{lemma}[Changing a local representative]
  \label{lem:picRankOnePresentation_change}
  \provenancetag{CUSTOM}
  \uses{def:picRankOneLocalPresentation}
  \notready
  Suppose \(\mathcal L\) and \(\mathcal L'\) represent the same relative Picard class over a
  test \(T\).  After an etale refinement, write
  \(\mathcal L'\simeq\mathcal L\otimes f_T^*\mathcal M\) for an invertible sheaf
  \(\mathcal M\) on \(T\).  Then \(\mathcal L\) satisfies the rank-one conditions of
  Definition~\ref{def:picRankOneLocalPresentation} if and only if \(\mathcal L'\) does.
  Under this equivalence their evaluation morphisms determine the same closed subscheme of
  \(C_T\).
\end{lemma}
\begin{proof}
  The projection formula gives functorial isomorphisms
  \[
    f_{T*}\mathcal L'\simeq f_{T*}\mathcal L\otimes\mathcal M,
    \qquad
    R^1f_{T*}\mathcal L'\simeq R^1f_{T*}\mathcal L\otimes\mathcal M.
  \]
  Tensoring by the invertible sheaf \(\mathcal M\) preserves vanishing, invertibility, and
  all base-change isomorphisms.  Under the first displayed isomorphism, the counit
  \(f_T^*f_{T*}\mathcal L'\to\mathcal L'\) is the counit for \(\mathcal L\) tensored by
  \(f_T^*\mathcal M\).  After tensoring source and target by the inverse line bundle, the two
  counits therefore have the same vanishing ideal, so their zero schemes agree.
\end{proof}

\begin{definition}[The rank-one Picard subfunctor]
  \label{def:picRankOneOpen}
  \provenancetag{CUSTOM}
  \uses{def:picRankOneLocalPresentation,
    lem:picRankOnePresentation_change}
  \notready
  Define \(\mathrm{PicRankOneOpen}(C,T)\) to be the set of
  \(\lambda\in\mathrm{pic}^{g(C)}(C,T)\) admitting a rank-one local presentation.
  Pullback of the presenting cover and line bundles makes these sets a subfunctor
  \[
    \mathrm{PicRankOneOpen}(C)\subseteq\mathrm{pic}^{g(C)}(C).
  \]
\end{definition}

\begin{theorem}[Openness of the rank-one Picard locus]
  \label{thm:picRankOneOpen_isOpen}
  \provenancetag{ADAPTED}
  \uses{def:picRankOneOpen, lem:picRankOnePresentation_change,
    thm:h0_eq_deg_add_chi, thm:chi_moduleKSheaf_curve}
  \notready
  The inclusion
  \(\mathrm{PicRankOneOpen}(C)\hookrightarrow\mathrm{pic}^{g(C)}(C)\) is relatively
  representable by open immersions, and its formation commutes with arbitrary base change.
  \source{kleiman-picard}
\end{theorem}
\begin{proof}
  Pull a test class back to an etale cover on which it is represented by an invertible sheaf
  \(\mathcal L\).  Upper semicontinuity makes the locus where
  \(H^1(C_t,\mathcal L_t)=0\) open.  On that locus, cohomology and base change identifies
  \(f_*\mathcal L\otimes k(t)\) with \(H^0(C_t,\mathcal L_t)\), and Riemann--Roch at degree
  \(g(C)\) gives
  \[
    \dim H^0(C_t,\mathcal L_t)=g(C)+1-g(C)=1.
  \]
  Hence \(f_*\mathcal L\) is invertible there and its formation commutes with every base
  change.  Conversely the rank-one conditions include the same \(H^1\)-vanishing, so this
  open is exactly the desired locus.  Lemma~\ref{lem:picRankOnePresentation_change} shows
  that the opens obtained from different representatives agree on a common etale refinement;
  open subsets descend uniquely along an etale cover.  The construction by fibrewise
  vanishing and base change also shows compatibility with an arbitrary pullback of the test.
\end{proof}

\subsection{The canonical divisor and the Abel isomorphism}

\begin{definition}[The divisor rank-one open]
  \label{def:divRankOneOpen}
  \provenancetag{CUSTOM}
  \uses{def:genusDivisorRepresentation, def:genusAbelMap,
    def:picRankOneOpen, thm:picRankOneOpen_isOpen}
  \notready
  Let \(W\subseteq D_g\) be the open subscheme representing the inverse image of
  \(\mathrm{PicRankOneOpen}(C)\) under \(a_g\).  Thus a \(T\)-point of \(W\) is precisely a
  degree-\(g(C)\) relative effective divisor whose associated Picard class belongs to the
  rank-one locus.  Write \(\mathrm{DivRankOneOpen}(C)=W\).
\end{definition}
\begin{proof}
  The inverse image of a relatively representable open subfunctor along the represented Abel
  map is represented by the corresponding open subscheme of \(D_g\).  The pointwise
  description is the defining property of this fibre product.
\end{proof}

\begin{definition}[The restricted rank-one Abel map]
  \label{def:rankOneAbel}
  \provenancetag{CUSTOM}
  \uses{def:divRankOneOpen, def:genusAbelMap}
  \notready
  Restriction of \(a_g\) to \(W\) gives a natural transformation
  \[
    \operatorname{rankOneAbel}\colon
      h_W\longrightarrow\mathrm{PicRankOneOpen}(C).
  \]
\end{definition}

\begin{definition}[The canonical divisor of a rank-one class]
  \label{def:divisorOfRankOne}
  \provenancetag{ADAPTED}
  \uses{def:picRankOneLocalPresentation,
    lem:picRankOnePresentation_change, def:divRankOneOpen}
  \notready
  There is a natural transformation
  \[
    \operatorname{divisorOfRankOne}\colon
      \mathrm{PicRankOneOpen}(C)\longrightarrow h_W
  \]
  defined by the zero scheme of the canonical evaluation morphism.  Its value at
  \(\lambda\) is finite locally free of degree \(g(C)\), commutes with arbitrary base change,
  and has relative Picard class \(\lambda\).
  \source{kleiman-picard}
\end{definition}
\begin{proof}
  On a member \(T_i\) of a rank-one presentation, set
  \(\mathcal N_i=f_{i*}\mathcal L_i\).  The counit
  \(f_i^*\mathcal N_i\to\mathcal L_i\) is equivalently a canonical section
  \[
    s_i\in H^0\bigl(C_{T_i},
      \mathcal L_i\otimes f_i^*\mathcal N_i^{-1}\bigr).
  \]
  On every geometric fibre, base change identifies \(s_i\) with the nonzero generator of the
  one-dimensional vector space \(H^0(C_t,\mathcal L_{i,t})\).  Since the fibre curve is
  integral, this section is regular.  The fibrewise criterion for a relative effective
  Cartier divisor therefore shows that its zero scheme \(E_i\) is a relative effective
  Cartier divisor.  Its fibre degree is the degree of \(\mathcal L_{i,t}\), namely \(g(C)\);
  hence \(E_i\to T_i\) is proper and quasi-finite, therefore finite, and its flat fibres have
  constant length \(g(C)\), so it is finite locally free of that rank.

  On a common refinement, two representatives differ by a pullback from the base.  The
  projection-formula comparison in Lemma~\ref{lem:picRankOnePresentation_change} identifies
  their evaluation sections and hence their zero schemes.  Effective descent for closed
  subschemes glues the \(E_i\) uniquely to a divisor \(E\subset C_T\).  The same comparison
  after an arbitrary map \(T'\to T\) proves that \(E\) commutes with base change.  Finally,
  \[
    \mathcal O_{C_T}(E)|_{C_{T_i}}
      \simeq \mathcal L_i\otimes f_i^*\mathcal N_i^{-1},
  \]
  so \(\mathcal O(E)\) and \(\mathcal L_i\) differ only by a line bundle pulled back from the
  base.  Their relative Picard classes are equal to \(\lambda\), which also shows that \(E\)
  is a \(T\)-point of \(W\).  These constructions are canonical, and therefore assemble into
  the asserted natural transformation.
\end{proof}

\begin{theorem}[The two rank-one constructions are inverse]
  \label{thm:rankOneAbel_inverse}
  \provenancetag{CUSTOM}
  \uses{def:rankOneAbel, def:divisorOfRankOne}
  \notready
  The natural transformations \(\operatorname{rankOneAbel}\) and
  \(\operatorname{divisorOfRankOne}\) are mutually inverse.
\end{theorem}
\begin{proof}
  Starting with a rank-one class \(\lambda\), the construction of
  Definition~\ref{def:divisorOfRankOne} gives a divisor \(E\) with
  \([\mathcal O(E)]=\lambda\) in the relative Picard functor.  Applying the Abel map therefore
  returns \(\lambda\).

  Conversely, start with a divisor \(E\in W(T)\) and put
  \(\mathcal L=\mathcal O(E)\).  The canonical divisor section gives a morphism
  \(\mathcal O_T\to f_{T*}\mathcal L\).  Both sides are invertible, and on each geometric
  fibre this morphism is the nonzero element represented by the divisor \(E_t\); it is
  therefore an isomorphism.  Under this isomorphism the evaluation section is exactly the
  canonical section of \(\mathcal O(E)\), whose zero scheme is \(E\).  Thus applying
  \(\operatorname{divisorOfRankOne}\) after the Abel map returns \(E\).  Both equalities are
  compatible with pullback, so they are equalities of natural transformations.
\end{proof}

\begin{definition}[The rank-one Abel isomorphism]
  \label{def:rankOneAbelIso}
  \lean{AlgebraicGeometry.canonicalRankOneAbelIso}
  \provenancetag{CUSTOM}
  \uses{thm:rankOneAbel_inverse}
  \leanok
  The mutually inverse transformations form a natural isomorphism
  \[
    \operatorname{rankOneAbelIso}\colon
      h_W\xrightarrow{\ \sim\ }\mathrm{PicRankOneOpen}(C).
  \]
\end{definition}
\begin{proof}
  \leanok
  Use \(\operatorname{rankOneAbel}\) as the forward map and
  \(\operatorname{divisorOfRankOne}\) as the inverse; the two triangle identities are
  Theorem~\ref{thm:rankOneAbel_inverse}.
\end{proof}

\begin{theorem}[The rank-one Abel map is an open immersion]
  \label{thm:rankOneAbel_isOpenImmersion}
  \provenancetag{CUSTOM}
  \uses{def:rankOneAbelIso, thm:picRankOneOpen_isOpen}
  \notready
  The composite
  \[
    h_W\xrightarrow{\operatorname{rankOneAbelIso}}
    \mathrm{PicRankOneOpen}(C)\hookrightarrow\mathrm{pic}^{g(C)}(C)
  \]
  is relatively representable by open immersions.
\end{theorem}
\begin{proof}
  A natural isomorphism is relatively representable by isomorphisms, and the second arrow is
  relatively representable by open immersions by
  Theorem~\ref{thm:picRankOneOpen_isOpen}.  The composite of an isomorphism with an open
  immersion is an open immersion, and the same argument holds after every base change.
\end{proof}

\subsection{The separably closed cover}

\begin{definition}[Translated rank-one opens in degree zero]
  \label{def:pic0TranslatedRankOneOpen}
  \provenancetag{CUSTOM}
  \uses{def:picRankOneOpen, def:pic0Functor}
  \notready
  Let \(\mathcal M\) be an invertible sheaf of degree \(g(C)\) on \(C\).  Define
  \[
    U_{\mathcal M}(T)=
      \{\lambda\in\mathrm{pic}^0(C,T):
        \lambda\otimes\mathcal M_T\in\mathrm{PicRankOneOpen}(C,T)\}.
  \]
  Translation by \(\mathcal M\) identifies \(U_{\mathcal M}\) with
  \(\mathrm{PicRankOneOpen}(C)\).
\end{definition}

\begin{theorem}[Representability of each translated open]
  \label{thm:pic0TranslatedRankOneOpen_representable}
  \provenancetag{CUSTOM}
  \uses{def:pic0TranslatedRankOneOpen, def:rankOneAbelIso,
    thm:rankOneAbel_isOpenImmersion}
  \notready
  Each \(U_{\mathcal M}\) is represented by a scheme, and its inclusion in
  \(\mathrm{pic}^0(C)\) is relatively representable by open immersions.
\end{theorem}
\begin{proof}
  Tensoring by \(\mathcal M_T\) is a natural isomorphism from the degree-zero layer to the
  degree-\(g(C)\) layer, with inverse tensoring by \(\mathcal M_T^{-1}\).  It carries
  \(U_{\mathcal M}\) isomorphically onto \(\mathrm{PicRankOneOpen}(C)\), which is represented
  by \(W\) through Definition~\ref{def:rankOneAbelIso}.  Transporting that representation
  and the open-immersion certificate through the translation proves both assertions.
\end{proof}

\begin{theorem}[Translated rank-one opens cover over a separably closed field]
  \label{thm:rankOne_translate_cover_sepClosed}
  \provenancetag{CUSTOM}
  \uses{def:pic0TranslatedRankOneOpen, thm:exists_rationalPoint_mem,
    thm:h0_eq_deg_add_chi, thm:chi_moduleKSheaf_curve}
  \notready
  Suppose \(k\) is separably closed.  For every field extension \(K/k\) and every
  \(\lambda\in\mathrm{pic}^0(C,\Spec K)\), there is an invertible sheaf
  \(\mathcal M\) of degree \(g(C)\) on \(C\) such that
  \(\lambda\in U_{\mathcal M}(\Spec K)\).
\end{theorem}
\begin{proof}
  Membership in the rank-one locus is etale local, so after a finite separable extension of
  \(K\) we may represent \(\lambda\) by an invertible sheaf \(\mathcal L\) on the base-changed
  curve.  Fix a positive invertible sheaf \(\mathcal A\) on \(C\).  Serre vanishing gives an
  \(n\) such that
  \(H^1(\mathcal L\otimes\mathcal A_K^n)=0\), and increasing \(n\) if necessary makes its
  space of sections nonzero.  Write this sheaf as \(\mathcal L_0\).

  If \(h^0(\mathcal L_j)>1\), choose a nonzero section \(s_j\).  Its zero scheme has finite
  support.  The \(k\)-rational points of \(C\) are infinite and dense by
  Theorem~\ref{thm:exists_rationalPoint_mem}, so there is a point \(P_j\in C(k)\) outside that
  support.  Evaluation at \(P_j\) is then onto, and the exact sequence
  \[
    0\longrightarrow\mathcal L_j(-P_j)\longrightarrow\mathcal L_j
      \longrightarrow\mathcal L_j|_{P_j}\longrightarrow0
  \]
  shows that \(h^0\) drops by one while \(H^1\) remains zero.  Repeating this finite process
  gives
  \[
    \mathcal L_r=\mathcal L\otimes
      \bigl(\mathcal A^n(-P_0-\cdots-P_{r-1})\bigr)_K
  \]
  with \(h^0(\mathcal L_r)=1\) and \(H^1(\mathcal L_r)=0\).  Put
  \(\mathcal M=\mathcal A^n(-P_0-\cdots-P_{r-1})\), which is defined over \(k\).
  Riemann--Roch gives
  \(1=\deg(\mathcal L_r)+1-g(C)\).  Since \(\lambda\) has degree zero, this says
  \(\deg\mathcal M=g(C)\).  Cohomology and base change now gives a rank-one local
  presentation of \(\lambda\otimes\mathcal M_K\).  The construction was made after an etale
  extension of \(K\), so etale descent of membership places the original class in
  \(U_{\mathcal M}(\Spec K)\).
\end{proof}

\begin{definition}[The separably closed degree-zero Picard representer]
  \label{def:pic0_sepClosed_representableBy}
  \lean{AlgebraicGeometry.pic0_sepClosed_representableBy}
  \provenancetag{CUSTOM}
  \uses{thm:pic0TranslatedRankOneOpen_representable,
    thm:rankOne_translate_cover_sepClosed, def:pic0RepresentableByOfCharts}
  \leanok
  If \(k\) is separably closed, there is a scheme \(J_s\) over \(k\) and a representation
  \[
    h_{J_s}\xrightarrow{\ \sim\ }\mathrm{pic}^0(C).
  \]
\end{definition}
\begin{proof}
  \leanok
  Use the represented opens \(U_{\mathcal M}\), indexed by the degree-\(g(C)\) invertible
  sheaves on \(C\).  Let \(T\) be a test scheme, \(\lambda\in\mathrm{pic}^0(C,T)\), and
  \(t\in T\).  Apply Theorem~\ref{thm:rankOne_translate_cover_sepClosed} to the residue field
  \(k(t)\) and the restriction \(\lambda_t\).  It supplies an \(\mathcal M\) for which \(t\)
  belongs to the open inverse image of \(U_{\mathcal M}\) in \(T\).  As \(t\) varies, these
  inverse images form an open cover of \(T\).  Hence the represented open immersions are
  jointly Zariski-locally surjective.  The local representability theorem
  (Definition~\ref{def:pic0RepresentableByOfCharts}) glues them and gives the asserted
  representation.
\end{proof}

\subsection{Finite Galois descent and the Jacobian datum}

\begin{definition}[The separably closed descent package]
  \label{def:pic0SepClosedDescentInput}
  \provenancetag{CUSTOM}
  \uses{def:pic0Functor}
  \notready
  A separably closed descent package consists of a scheme \(J_s\) representing
  \(\mathrm{pic}^0(C)\), its universal Picard class, and certificates that \(J_s\) is of
  finite type, proper, geometrically irreducible, and a group scheme.  The representation and
  the universal class in the package are the same datum, related by Yoneda, rather than two
  independently chosen witnesses.
\end{definition}

\begin{theorem}[The rank-one representer supplies the descent package]
  \label{thm:pic0SepClosedDescentInput_exists}
  \provenancetag{CUSTOM}
  \uses{def:pic0_sepClosed_representableBy,
    thm:genusDivisorSource_geometry, def:genusAbelMap,
    def:pic0SepClosedDescentInput}
  \notready
  Over a separably closed field, the representation of
  Definition~\ref{def:pic0_sepClosed_representableBy} extends to a separably closed descent
  package.
\end{theorem}
\begin{proof}
  The commutative-group law on \(\mathrm{pic}^0(C)\) transfers uniquely across the
  representation to a group-scheme law on \(J_s\), and the image of the identity morphism is
  its universal class.  The translated rank-one charts show that \(J_s\) is locally of finite
  type.  Its identity point is closed because it is a rational point of a scheme locally of
  finite type over a field.  The diagonal is the inverse image of that point under the
  difference morphism \(J_s\times J_s\to J_s\), so the group scheme is separated.

  Choose a point of \(C\) over the separably closed base and use its \(g(C)\)-fold divisor to
  translate the degree-\(g(C)\) Abel map to \(\mathrm{pic}^0(C)\).  After any algebraically
  closed extension, Riemann--Roch gives a section of every degree-\(g(C)\) line bundle, so this
  translated Abel map \(D_g\to J_s\) is surjective.  It is proper because \(D_g\) is projective
  and \(J_s\) is separated.  A surjective image of the quasi-compact scheme \(D_g\) makes
  \(J_s\) quasi-compact; combined with local finite type, this gives finite type.  Properness
  descends through the surjective proper map: after every base change, the image in the base
  of a closed subset of \(J_s\) is the image of its closed inverse image in \(D_g\), hence is
  closed.  Finally, the geometric fibres of \(D_g\) are irreducible by
  Theorem~\ref{thm:genusDivisorSource_geometry}, and their surjective images are irreducible.
  These are exactly the certificates required by the package.
\end{proof}

\begin{theorem}[Limit and faithfully flat compatibility of \(\mathrm{pic}^0\)]
  \label{thm:pic0DescentCompatibility}
  \provenancetag{CUSTOM}
  \uses{def:pic0Functor}
  \notready
  The degree-zero etale Picard functor commutes with filtered unions of field extensions and
  satisfies effective descent along faithfully flat field extensions, compatibly with its
  universal classes.
\end{theorem}
\begin{proof}
  An invertible sheaf is given on a finite affine cover by finitely many finitely presented
  rank-one modules, transition isomorphisms, and finitely many cocycle equations.  All this
  data, as well as an isomorphism between two such data, descends to a finite stage of a
  filtered union.  The degree-zero condition is preserved and reflected after field extension
  because degree is invariant under scalar extension.  This proves filtered compatibility.
  Faithfully flat descent for finitely presented modules descends the local line bundles and
  their transition maps, while the cocycle equation gives effectivity and uniqueness.  Passing
  to the quotient by line bundles from the base and then to the etale sheafification preserves
  this descent statement, since both equivalence and effectivity may be checked after a common
  faithfully flat refinement.
\end{proof}

\begin{theorem}[A finite Galois stage for the Picard representer]
  \label{thm:pic0FiniteGaloisStage}
  \provenancetag{CUSTOM}
  \uses{def:pic0SepClosedDescentInput, thm:pic0SepClosedDescentInput_exists,
    thm:pic0DescentCompatibility}
  \notready
  Let \(k^s\) be a separable closure of \(k\).  The separably closed descent package descends
  to a finite Galois extension \(L/k\): there are a finite-type proper group scheme \(J_L\)
  representing \(\mathrm{pic}^0(C_L)\), its universal class, and a semilinear action of
  \(\operatorname{Gal}(L/k)\) preserving that class and satisfying the cocycle condition.
  Every orbit of this action is contained in an affine open of \(J_L\).
\end{theorem}
\begin{proof}
  Because \(J_{k^s}\) is of finite presentation, it descends to a finite subextension
  \(L_0/k\).  The universal class descends after enlarging \(L_0\), by
  Theorem~\ref{thm:pic0DescentCompatibility}.  The natural transformation defined by the
  descended class becomes the given representation after extension to \(k^s\); faithful flat
  descent for both the Yoneda functor and the Picard sheaf shows that it is already an
  isomorphism over \(L_0\).  Properness, the group operations, and their finitely many axioms
  also descend after one further finite extension.  Enlarge \(L_0\) to a finite Galois
  extension \(L/k\).

  For \(\sigma\in\operatorname{Gal}(L/k)\), both \(\sigma^*J_L\) and \(J_L\), equipped with
  their transported universal classes, represent the same functor.  Yoneda uniqueness gives
  a unique semilinear isomorphism between them.  The two composites associated to
  \(\sigma\tau\) carry the universal class to the same element, so uniqueness proves the
  cocycle equation.  A proper finite-type group scheme over a field is projective.  Given a
  finite Galois orbit, choose a closed specialization of each of its finitely many points.
  For a sufficiently high power of an ample line bundle, restriction of global sections to
  the resulting finite closed subscheme is surjective.  Lift the section that is \(1\) at
  every member of that subscheme.  Its nonvanishing locus contains the chosen closed
  specializations and therefore the original orbit; as the complement of an ample effective
  divisor, this locus is affine.  This proves the final certificate.
\end{proof}

\begin{lemma}[Stable affine neighbourhoods for a finite action]
  \label{lem:finiteGalois_stableAffineCover}
  \provenancetag{CUSTOM}
  \notready
  Let a finite group \(G\) act on a separated scheme \(X\).  If every \(G\)-orbit is contained
  in an affine open, then \(X\) has a \(G\)-stable affine open cover.
\end{lemma}
\begin{proof}
  Fix \(x\in X\) and choose an affine open \(V\) containing its whole orbit.  Every translate
  \(gV\) also contains that orbit.  Hence
  \(U_x=\bigcap_{g\in G}gV\) is a \(G\)-stable open containing \(x\).  A finite intersection
  of affine opens in a separated scheme is affine: the intersection of two is the inverse
  image of the affine \(V\times V'\) under the affine diagonal, and induction handles the
  finite intersection.  The opens \(U_x\), as \(x\) varies, form the required cover.
\end{proof}

\begin{theorem}[Finite Galois descent of a scheme]
  \label{thm:finiteGalois_schemeDescent}
  \provenancetag{CUSTOM}
  \uses{lem:finiteGalois_stableAffineCover}
  \notready
  Let \(L/k\) be finite Galois and let an \(L\)-scheme \(X_L\) carry a semilinear Galois
  action satisfying the cocycle condition.  If every orbit lies in an affine open, then there
  are a \(k\)-scheme \(X\) and an isomorphism
  \(X\times_kL\simeq X_L\) inducing the given action.
\end{theorem}
\begin{proof}
  By Lemma~\ref{lem:finiteGalois_stableAffineCover}, cover \(X_L\) by stable affine opens
  \(U_i=\Spec A_i\).  Semilinear Galois descent for algebras gives
  \(B_i=A_i^G\) and an isomorphism \(L\otimes_kB_i\simeq A_i\).  Since \(X_L\) is separated,
  every overlap \(U_i\cap U_j\) and every triple overlap is affine and stable.  Applying the
  same invariant construction to these overlaps descends both open immersions
  \(U_i\cap U_j\rightrightarrows U_i,U_j\); faithful flatness of \(L/k\) reflects the open
  immersions and their cocycle identities.  The affine schemes \(\Spec B_i\) therefore glue
  along the descended overlaps to a \(k\)-scheme \(X\).  The local isomorphisms
  \(L\otimes_kB_i\simeq A_i\) agree on overlaps and glue to
  \(X\times_kL\simeq X_L\).  Their defining invariant maps identify the canonical descent
  action with the prescribed semilinear action.
\end{proof}

\begin{theorem}[Finite Galois descent of a representation]
  \label{thm:representableBy_of_finiteGalois_baseChange}
  \provenancetag{CUSTOM}
  \uses{thm:finiteGalois_schemeDescent}
  \notready
  Let \(F\) be a set-valued sheaf on \(k\)-schemes satisfying faithfully flat descent, and let
  \(L/k\) be finite Galois.  Suppose \(F_L\) is represented by \(X_L\), and the universal
  element is preserved by a semilinear Galois action on \(X_L\) satisfying the cocycle
  condition and the orbit-in-affine hypothesis.  Then there are a \(k\)-scheme \(X\) and a
  representation \(h_X\simeq F\); the pair \((X,h_X\simeq F)\) is returned together, and its
  base change is the given pair over \(L\).
\end{theorem}
\begin{proof}
  Descend \(X_L\) to \(X\) by Theorem~\ref{thm:finiteGalois_schemeDescent}.  The invariance of
  the universal element is precisely a descent datum on its two pullbacks to
  \(X_L\times_X X_L\).  The sheaf condition gives a unique element \(u\in F(X)\) descending
  it.  Yoneda turns \(u\) into a natural map \(h_X\to F\).  After base change to \(L\) this is
  the given representing isomorphism.  For any test \(T\), apply the sheaf equalizer diagrams
  for the finite faithfully flat cover \(T_L\to T\) to both \(h_X\) and \(F\).  The maps on
  \(T_L\) and \(T_L\times_TT_L\) are bijections and commute with the two restriction maps;
  hence the map on their equalizers is a bijection.  Thus \(u\) represents \(F\).  The
  construction keeps \(X\) and this certificate in the same pair, and the glued base-change
  isomorphism identifies it with the original represented pair over \(L\).
\end{proof}

\begin{definition}[The arbitrary-field degree-zero Picard representer]
  \label{def:pic0_representableBy}
  \provenancetag{CUSTOM}
  \uses{thm:pic0FiniteGaloisStage,
    thm:representableBy_of_finiteGalois_baseChange}
  \notready
  For every field \(k\), there is a distinguished pair
  \[
    \bigl(J,\;h_J\xrightarrow{\ \sim\ }\mathrm{pic}^0(C)\bigr)
  \]
  obtained by finite Galois descent from the separably closed rank-one construction.
\end{definition}
\begin{proof}
  Apply Theorem~\ref{thm:pic0FiniteGaloisStage} to obtain the represented finite Galois model,
  its action, cocycle, invariant universal class, and orbit-in-affine certificate.  The Picard
  functor has the descent property of Theorem~\ref{thm:pic0DescentCompatibility}, so
  Theorem~\ref{thm:representableBy_of_finiteGalois_baseChange} returns the required scheme and
  representation certificate as one pair.  This is the asserted distinguished datum.
\end{proof}

\begin{definition}[The Jacobian datum from the chosen representation]
  \label{def:jacobianData}
  \provenancetag{CUSTOM}
  \uses{def:pic0_representableBy, def:jacobian_data,
    thm:pic0FiniteGaloisStage}
  \notready
  Let \((J,r)\) be the pair of
  Definition~\ref{def:pic0_representableBy}.  The \emph{rank-one descent Jacobian datum} is
  the Jacobian datum whose scheme is \(J\), whose representation is exactly \(r\), and whose
  finite-type and quasi-compact certificates are the descents of the corresponding
  certificates at the finite Galois stage.
\end{definition}
\begin{proof}
  Finite type and properness descend and are reflected by the finite faithfully flat extension
  \(L/k\) in Theorem~\ref{thm:pic0FiniteGaloisStage}; in particular \(J\to\Spec k\) is locally
  of finite type and quasi-compact.  Insert \(J\), the same certificate \(r\), and
  these two properties into Definition~\ref{def:jacobian_data}.  No second representing object
  or universal element is chosen: the group law and universal class are transported from the
  same \(r\) by Yoneda.
\end{proof}
